\chapter{Solar System \& Planetary Science}

\section*{Theory problems}

\probhead{3.1}{}{2007}{Chiang Mai, Thailand}{TH 1.2}{2}
% source id: 2007_TH_1.2   tag: Venus greatest elongation gives orbit radius
The largest angular separation between Venus and the Sun, when viewed from the
Earth, is $46^\circ$. Calculate the radius of Venus's circular orbit in A.U.

\probhead{3.2}{}{2007}{Chiang Mai, Thailand}{TH 1.4}{2}
% source id: 2007_TH_1.4   tag: Moon distance from angular diameter
One night during a full Moon, the Moon subtends an angle of 0.46 degree to an
observer. What is the observer's distance to the Moon on that night?

\probhead{3.3}{A Planet \& Its Surface Temperature}{2007}{Chiang Mai, Thailand}{TH Q2}{10}
% source id: 2007_TH_Q2   tag: planet equilibrium temperature, albedo, rotation
A fast rotating planet of radius $R$ with surface albedo $a$ is orbiting a star
of luminosity $L$. The orbital radius is $D$. It is assumed here that, at
equilibrium, all of the energy absorbed by the planet is re-emitted as a
blackbody radiation.

\begin{description}
\item[a.)] What is the radiation flux from the star at the planet's surface?
  \hfill(1.5 points)
\item[b.)] What is the total rate of energy absorbed by the planet?
  \hfill(1.5 points)
\item[c.)] What is the reflected luminosity of the planet? \hfill(2 points)
\item[d.)] What is the average blackbody temperature of the planet's surface?
  \hfill(2 points)
\item[e.)] If we were to assume that one side of the planet is always facing
  the star, what would be the average surface temperature of that side?
  \hfill(2 points)
\item[f.)] For the planet in problem d:
  \[ a = 0.25, \qquad D = 1.523\ \mathrm{A.U.}, \]
  calculate its surface temperature in kelvins for the value of
  $L = 3.826 \times 10^{26}$\,W. \hfill(1 point)
\end{description}

\probhead{3.4}{}{2008}{Bandung, Indonesia}{TH S11}{20}
% source id: 2008_TH_S11   tag: annular eclipse photo: distances, coverage
Below is a picture on a 35\,mm film of annular solar eclipse in Dumai, Riau,
Indonesia on August 22, 1998, taken with a telescope having effective diameter
10\,cm and f-ratio 15. The diameter of the Sun's disk in original picture on
the film is 13.817\,mm and the diameter of the Moon's disk is 13.235\,mm.
Determine the distances of the Sun and the Moon (expressed in km) from the
Earth and the percentage of the solar disk covered by the Moon during the
annular solar eclipse.
\ioaafig{2008_TH_S11-f1.pdf}
% disk with the dark lunar disk inside), printed below the statement in the
% 2008 theory paper, problem 11.

\probhead{3.5}{}{2008}{Bandung, Indonesia}{TH S13}{20}
% source id: 2008_TH_S13   tag: tangent-line angles to circular asteroid
In the journey of a space craft, scientists make a close encounter with an
object and they would like to investigate the object more carefully using their
on-board telescope. For simplicity, we assume this to be a two-dimensional
problem and that the position of the space craft is stationary in $(0,0)$. The
shape of the object is a disk and the boundary has the equation
\[ x^2 + y^2 - 10x - 8y + 40 = 0. \]
Find the exact values of maximum and minimum of $\tan\varphi$ where $\varphi$
is the angle of the telescope with respect to the $x$ direction during
investigation from one edge to the other edge.

\probhead{3.6}{}{2008}{Bandung, Indonesia}{TH S9}{20}
% source id: 2008_TH_S9   tag: Earth temperature with atmospheric transmittance
The radiation incoming to the Earth from the Sun must penetrate the Earth's
atmosphere before reaching the earth surface. The Earth also releases radiation
to its environment and this radiation must penetrate the Earth's atmosphere
before going out to the outer space. In general, the transmittance ($t_1$) of
the Sun radiation during its penetration through the Earth's atmosphere is
higher than that of the radiation from the Earth ($t_2$). Let
$T_{\mathrm{eff}\,\odot}$ be the effective temperature of the Sun, $R_\odot$
the radius of the Sun, $r_\oplus$ the radius of the Earth, and $x$ the distance
between the Sun and the Earth. Derive the temperature of the Earth's surface as
a function of the aforementioned parameters.

\probhead{3.7}{}{2009}{Tehran, Iran}{TH P13}{10}
% source id: 2009_TH_P13   tag: Mars-Sun distance, opposition-quadrature timing
Assume that you are living in the time of Copernicus and do not know anything
about Kepler's laws. You might calculate Mars--Sun distance in the same way as
he did. After accepting the revolutionary belief that all the planets are
orbiting around the Sun, not around the Earth, you measure that the orbital
period of Mars is 687 days, then you observe that 106 days after opposition of
Mars, the planet appears in quadrature. Calculate Mars--Sun distance in
astronomical unit (AU).

\probhead{3.8}{}{2010}{Beijing, China}{TH L17}{30}
% source id: 2010_TH_L17   tag: fictional planetary system days, moons
The planet Taris is the home of the Korribian civilization. The Korribian
species is a highly intelligent alien life form. They speak Korribianese
language. The Korribianese--English dictionary is shown in Table~1; read it
carefully! Korriban astronomers have been studying the heavens for thousands of
years. Their knowledge can be summarized as follows:

\begin{itemize}
\item Taris orbits its host star Sola in a circular orbit, at a distance of
  1 Tarislength.
\item Taris orbits Sola in 1 Tarisyear.
\item The inclination of Taris's equator to its orbital plane is $3^\circ$.
\item There are exactly 10 Tarisdays in 1 Tarisyear.
\item Taris has two moons, named Endor and Extor. Both have circular orbits.
\item The sidereal orbital period of Endor (around Taris) is exactly
  0.2 Tarisdays.
\item The sidereal orbital period of Extor (around Taris) is exactly
  1.6 Tarisdays.
\item The distance between Taris and Endor is 1 Endorlength.
\item Corulus, another planet, also orbits Sola in a circular orbit. Corulus
  has one moon.
\item The distance between Sola and Corulus is 9 Tarislengths.
\item The tarisyear begins when Solaptic longitude of the Sola is zero.
\end{itemize}

\begin{center}
\begin{tabular}{ll}
\multicolumn{2}{l}{Table 1: Korribianese--English dictionary}\\[4pt]
\textbf{Korribianese} & \textbf{English Translation}\\
Corulus & A planet orbiting Sola\\
Endor & (i) Goddess of the night; (ii) a moon of Taris\\
Endorlength & The distance between Taris and Endor\\
Extor & (i) God of peace; (ii) a moon of Taris\\
Sola & (i) God of life; (ii) the star which Taris and Corulus orbit\\
Solaptic & Apparent path of Sola and Corulus as viewed from Taris\\
Taris & A planet orbiting the star Sola, home of the Korribians\\
Tarisday & The time between successive midnights on the planet Taris\\
Tarislength & The distance between Sola and Taris\\
Tarisyear & Time taken by Taris to make one revolution around Sola\\
\end{tabular}
\end{center}

Questions:
\begin{description}
\item[(a)] Draw the Sola-system, and indicate all planets and moons.
  \hfill(5 points)
\item[(b)] How often does Taris rotate around its axis during one Tarisyear?
  \hfill(5 points)
\item[(c)] What is the distance between Taris and Extor, in Endorlengths?
  \hfill(3 points)
\item[(d)] What is the orbital period of Corulus, in Tarisyears?
  \hfill(2 points)
\item[(e)] What is the distance between Taris and Corulus when Corulus is in
  opposition? \hfill(5 points)
\item[(f)] If at the beginning of a particular tarisyear, Corulus and taris
  were in opposition, what would be Solaptic longitude (as observed from Taris)
  of Corulus $n$ tarisdays from the start of that year? \hfill(5 points)
\item[(g)] What would be the area of the triangle formed by Sola, Taris and
  Corulus exactly one tarisday after the opposition? \hfill(5 points)
\end{description}

\probhead{3.9}{}{2010}{Beijing, China}{TH S11}{10}
% source id: 2010_TH_S11   tag: next Mars great opposition, synodic period
Mars arrived at its great opposition at UT $17^{\mathrm{h}}56^{\mathrm{m}}$
Aug.\ 28, 2003. The next great opposition of Mars will be in 2018, estimate the
date of that opposition. The semi-major axis of the orbit of Mars is
1.524\,AU.

\probhead{3.10}{}{2010}{Beijing, China}{TH S14}{10}
% source id: 2010_TH_S14   tag: Venus transit chord geometry
An observer observed a transit of Venus near the North Pole of the Earth. The
transit path of Venus is shown in the picture below. A, B, C, D are all on the
path of transit and marking the center of the Venus disk. At A and B, the
center of Venus is superposed on the limb of the Sun disk; C corresponds to the
first contact while D to the fourth contact, $\angle AOB = 90^\circ$, MN is
parallel to AB. The first contact occured at 9:00 UT. Calculate the time of the
fourth contact.
\[ T_{\mathrm{venus}} = 224.70\ \mathrm{days}, \quad
   T_{\mathrm{earth}} = 365.25\ \mathrm{days}, \quad
   a_{\mathrm{venus}} = 0.723\,\mathrm{AU}, \quad
   r_{\mathrm{venus}} = 0.949\,r_\oplus \]
\ioaafig{2010_TH_S14-f1.pdf}
% through points A, B (centre of Venus on the solar limb), C (first contact),
% D (fourth contact), and line MN parallel to AB; printed in the 2010
% proceedings (solutions section, p. 38, left figure).

\probhead{3.11}{}{2010}{Beijing, China}{TH S15}{10}
% source id: 2010_TH_S15   tag: annular vs total eclipse ratio
On average, the visual diameter of the Moon is slightly less than that of the
Sun, so the frequency of annular solar eclipses is slightly higher than total
solar eclipses. For an observer on the Earth, the longest total solar eclipse
duration is about 7.5 minutes, and the longest annular eclipse duration is
about 12.5 minutes. Here, the longest duration is the time interval from the
second contact to the third contact. Suppose we count the occurrences of both
types of solar eclipses for a very long time, estimate the ratio of the
occurrences of annular solar eclipses and total solar eclipses. Assume the
orbit of the Earth to be circular and the eccentricity of the Moon's orbit is
0.0549. Count all hybrid eclipses as annular eclipses.

\probhead{3.12}{}{2011}{Katowice, Poland}{TH S10}{10}
% source id: 2011_TH_S10   tag: proton energy to penetrate magnetosphere
Estimate the minimum energy a proton would need to penetrate the Earth's
magnetosphere. Assume that the initial penetration is perpendicular to a belt
of constant magnetic field $30\,\mu\mathrm{T}$ and thickness
$1.0 \times 10^{4}$\,km. Prepare the sketch of the particle trajectory. (Note
that at such high energies the momentum can be replaced by the expression
$E/c$. Ignore any radiative effects).

\probhead{3.13}{}{2012}{Rio de Janeiro, Brazil}{TH S3}{}
% source id: 2012_TH_S3   tag: synodic period between Mars oppositions
What is the time interval between two consecutive oppositions of Mars? Assume
the orbit is circular.

\probhead{3.14}{}{2012}{Rio de Janeiro, Brazil}{TH S4}{}
% source id: 2012_TH_S4   tag: full-Moon magnitude if albedo unity
What would be full Moon's visual magnitude if its albedo were equal to 1?

\probhead{3.15}{}{2013}{Volos, Greece}{TH S1}{10}
% source id: 2013_TH_S1   tag: Earth equilibrium temperature with albedo
What would be the mean temperature on the Earth's surface if we ignore the
greenhouse effect, assume that the Earth is a perfect black body and take into
account its non-vanishing albedo? Assume that the Earth's orbit around the Sun
is circular.

\probhead{3.16}{Apparent magnitude of the Moon}{2014}{Suceava, Romania}{TH P14}{}
% source id: 2014_TH_P14   tag: Moon magnitude at phases
You know that the absolute magnitude of the Moon is
$M_L = 0.25^{\mathrm{m}}$. Calculate the values of the apparent magnitudes of
the Moon corresponding to the following Moon-phases: full-moon and the first
quarter. You know: the Moon--Earth distance --
$d_{\mathrm{LP}} = 385000$\,km, the Earth--Sun distance --
$d_{\mathrm{PS}} = 1$\,AU, the Moon--Sun distance,
$d_{\mathrm{LS}} = 1$\,AU.

\probhead{3.17}{}{2015}{Magelang, Indonesia}{TH S3}{}
% source id: 2015_TH_S3   tag: asteroid diameter from occultation
On 27 May 2015 at 02:18:49, the occultation of the star HIP 89931
($\delta$--Sgr) by the asteroid 1285 Julietta was observed from Borobudur
temple, which was located at the center of the asteroid shadow path. It lasted
for only 6.201\,s. Assume that Earth's orbit is circular and the orbit of
Julietta is on the ecliptic plane and revolves in the same direction as Earth.
At the occultation, Julietta is near its aphelion. At the time of the
occultation the distances of Julietta from the Sun and the Earth are
3.076\,AU and 2.156\,AU respectively. Find the approximate diameter of asteroid
Julietta, if the semi major axis of Julietta is $a = 2.9914$\,AU.

\probhead{3.18}{}{2015}{Magelang, Indonesia}{TH S9}{}
% source id: 2015_TH_S9   tag: Venus equilibrium temperature, radio flux
Assume that the Sun is a perfect blackbody. Venus is also assumed to be a
blackbody, with temperature $T_{\mathrm{V}}$, and it is in thermal equillibrium
(i.e.\ it is radiating about as much energy as it receives from the Sun) at its
orbital distance of 0.72\,AU. Suppose that at closest approach to Earth, Venus
has an angular diameter of about 66 arcsec. What is the flux density of Venus
at the closest approach to Earth as observed by a radio telescope at an
observing frequency of 5\,GHz?

\probhead{3.19}{Gases on Titan}{2016}{Bhubaneswar, India}{TH T2}{10}
% source id: 2016_TH_T2   tag: minimum atomic weight retained by Titan
Gas particles in a planetary atmosphere have a wide distribution of speeds. If
the r.m.s.\ (root mean square) thermal speed of particles of a particular gas
exceeds $1/6$ of the escape speed, then most of that gas will escape from the
planet. What is the minimum atomic weight (relative atomic mass),
$A_{\min}$, of an ideal monatomic gas so that it remains in the atmosphere of
Titan?

Given, mass of Titan $M_{\mathrm{T}} = 1.23 \times 10^{23}$\,kg, radius of
Titan $R_{\mathrm{T}} = 2575$\,km, surface temperature of Titan
$T_{\mathrm{T}} = 93.7$\,K.

\probhead{3.20}{The Star of Bethlehem}{2017}{Phuket, Thailand}{TH T8}{20}
% source id: 2017_TH_T8   tag: Jupiter-Saturn great conjunction period
A great conjunction is a conjunction of Jupiter and Saturn for observers on
Earth. Assume that Jupiter and Saturn have circular orbits in the ecliptic
plane.

The time between successive conjunctions may vary slightly as viewed from the
Earth. However, the average time period of the great conjunctions is the same
as that of an observer at the centre of the Solar system.

\begin{description}
\item[a)] Find the average great conjunction period (in years) and average
  heliocentric angle between two successive great conjunctions (in degrees).
  \hfill[6]
\item[b)] The next great conjunction will be on 21st December 2020 with an
  elongation of $30.3^\circ$ East of the Sun. Suggest the constellation in
  which the conjunction on 21st December 2020 will occur. (Give the IAU Latin
  name or IAU three-letter abbreviation of the constellation, i.e.\ Ursa Major
  or UMa) \hfill[2]
\end{description}

In 1606, Johannes Kepler determined that in some years the great conjunction
can be happen three times in the same year due to the retrograde motions of
% sic: "can be happen" as in the original
the planets. He also determined that such an event happened in the year 7 BC,
which could have been the event commonly known as ``The Star of Bethlehem''.
For the calculations below you may ignore the precession of the axis of the
Earth.

\begin{description}
\item[c)] Suggest the constellation in which the great conjunctions in 7 BC
  occurred. (Give the IAU Latin name or IAU three-letter abbreviation of the
  constellation, i.e.\ Ursa Major or UMa) \hfill[8]
\item[d)] During the second conjunction of the series of three conjunctions
  in 7 BC, suggest the constellation the Sun was in as viewed by the observer
  on Earth. (Give the IAU Latin name or IAU three-letter abbreviation of the
  constellation, i.e.\ Ursa Major or UMa) \hfill[4]
\end{description}

\probhead{3.21}{Famous astronomical events}{2019}{Keszthely, Hungary}{TH 1}{10}
% source id: 2019_TH_1   tag: chronological ordering of events
Arrange the following astronomy-related events in chronological order from the
oldest to the most recent. Write the correct serial number (between 1 and 11)
in the appropriate box on the answer sheet.

\begin{enumerate}
\item Launch of the Hubble Space Telescope
\item Viking probes arrived at planet Mars
\item Discovery of Phobos and Deimos
\item Latest perihelion of comet 1/P Halley
\item Discovery of Ceres (asteroid / dwarf planet)
\item Discovery of Uranus (planet)
\item First successful measurement of a stellar parallax
\item Discovery of the first planetary nebula
\item Discovery of stellar populations (I and II)
\item First identification of a quasar with an optical source
\item Discovery of the expansion of the Universe
\end{enumerate}

\probhead{3.22}{Flat Earth}{2020}{GeCAA (Estonia, remote)}{TH 2}{10}
% source id: 2020_TH_2   tag: flat-Earth equilibrium temperature
A new model of the world is gaining in popularity among some people. These
people believe in the ``Flat Earth'' view of the world, where the Earth is not
a spheroid, but rather a circle with radius $R_\oplus$. The central axis of
the Earth (normal to the circle passing through its centre C) is passes
through the observer's zenith.
% sic: "is passes through" as in the original

This model must at least remain consistent with the observed phenomena, as
listed below:
\begin{itemize}
\item The value of the solar constant is
  $S_\odot = 1366\,\mathrm{W/m}^2$.
\item The Earth's central axis precesses in a circle with a period
  25800 years.
\item The radius of the precession circle is $23.5^\circ$.
\end{itemize}

We assume that the Earth is a perfect blackbody radiator and the Sun is
sufficiently far away that all sun rays are parallel. Let us also assume that
the Sun's current (initial) location is at the zenith.

Determine how many years it will take for the Earth's equilibrium temperature
to decrease by $1\,^\circ$C.

\probhead{3.23}{Jupiter's Great Red Spot}{2020}{GeCAA (Estonia, remote)}{TH 8}{30}
% source id: 2020_TH_8   tag: GRS vorticity, spheroid curvature, Coriolis
In the following problem the fluid mechanics of Jupiter's Great Red Spot
(GRS) is studied based on the velocity field data. The diagram below shows a
map of relative velocity for GRS and the surrounding region. The arrows are
oriented and scaled as per the directions and magnitudes of winds at
different points.

\ioaafig{2020_TH_8-f1.pdf}

Due to the combined effects of gravity and rotation, Jupiter is slightly
flattened at its poles. The equation of a spheroid approximating for the
shape of Jupiter can be stated as:
\[ \frac{x^2 + y^2}{R_e^2} + \frac{z^2}{R_p^2} = 1, \]
where $R_e = 7.15 \times 10^7$\,m is the equatorial radius of Jupiter, and
$R_p = 6.69 \times 10^7$\,m the polar radius. The radii of curvature of this
spheroid in any direction can be calculated by the following equations
($\epsilon = \frac{R_e}{R_p}$):
\[ r(\phi) = R_e \left( 1 + \epsilon^{-2} \tan^2 \phi \right)^{-1/2} \]
\[ R(\phi) = R_e \epsilon^{-2}
   \left( \frac{r(\phi)}{R_e \cos\phi} \right)^{3} \]
where $r$ and $R$ are the zonal (aka in the zone of a particular latitude)
and meridional (aka longitudinal) radii of curvature, respectively, as a
function of planetographic latitude $\phi$. The sidereal rotation period of
Jupiter is $P = 3.57 \times 10^4$\,s.

\begin{description}
\item[(a)] (4 points) Calculate the zonal and meridional radii values ($\overline{r}$
  and $\overline{R}$ respectively) at the location of the centre of the GRS.
\item[(b)] (5 points) Estimate the eccentricity of the GRS.
\item[(c)] (6 points) 'Vorticity' at any point is a measure of local spinning
  of the fluid as measured by an observer situated in the reference frame of
  the fluid. Mathematically, it is calculated as 'curl' (vector derivative
  product) of the velocity field. In this case, the average relative
  vorticity may be estimated by the equation:
  \[ \xi = \frac{V_w L_{GRS}}{A_{GRS}} \]
  where $V_w$ is the maximum value of winds as per the velocity field,
  $L_{GRS}$ is the length of the circumference of the GRS and $A_{GRS}$ is
  the area of the GRS.

  Estimate average relative vorticity of the GRS.

  \textbf{Hint:} The circumference of an ellipse is well approximated by an
  average of circumferences of the corresponding auxiliary and minor circles.
\item[(d)] (2 points) Find the absolute vorticity $\xi_a = (\xi + f)$ by
  adding the Coriolis parameter
  \[ f = 2\Omega \sin\phi \]
  where $\Omega$ is the angular velocity of the Jupiter (due to axial
  rotation) and is the appropriate latitude.
  % sic: original omits "$\phi$" before "is the appropriate latitude"
\item[(e)] (1 point) If the absolute vorticity has the same sign as the
  latitude, we call the storm a `cyclonic storm'. If they have opposite
  signs, the system is `anticyclonic'. Is the GRS cyclonic or anticyclonic?
\item[(f)] (12 points) Imagine that the GRS moves to another latitude
  $\phi_1$, where the absolute vorticity changes the sign (changes from
  anti-cyclonic to cyclonic or vice versa). Assuming minimum possible
  displacement of the GRS, at what value of $\phi_1$ do we expect this
  change?

  In your analysis, assume that the GRS at the new location would occupy the
  same angular span in latitude, as well as have the same wind velocities and
  eccentricity as the original.
\end{description}

\probhead{3.24}{Temperature of the Earth}{2021}{Bogot\'a (remote)}{TH Q2}{10}
% source id: 2021_TH_Q2   tag: Earth equilibrium temperature, greenhouse
For at least the last few million years, the Earth has been in roughly
thermal equilibrium with the radiation from the Sun at the Earth's orbital
distance.

\begin{description}
\item[(2.1)] Assuming our planet to be an ideal black body, calculate what
  the Earth's equilibrium temperature (in Celsius) would be.
  \hfill[4.0\,pt]
\item[(2.2)] The Earth's albedo is approximately 30\%. Calculate the Earth's
  surface temperature (in Celsius) considering its albedo. \hfill[2.0\,pt]
\item[(2.3)] The Earth's absorbed radiation is reemitted as black body
  radiation from its surface, but its atmosphere re-absorbs 58\% of that
  energy, causing a greenhouse effect. Considering this effect, calculate the
  Earth's surface temperature (which will be the same as the temperature of
  the lower atmosphere). Give your answer in Celsius.

  For simplicity, consider the reabsorption effect as happening only once,
  and do not consider the atmosphere as a separate black body.
  \hfill[4.0\,pt]
\end{description}

\probhead{3.25}{Solar Retrograde Motion on Mercury}{2022}{Kutaisi, Georgia}{TH 9}{20}
% source id: 2022_TH_9   tag: Sun's retrograde motion in Mercury's sky
Mercury's orbit has an unusually high eccentricity. Further, its sidereal
rotational period is $\frac{2}{3}$ of its sidereal year. As a consequence of
these factors, the Sun will exhibit retrograde motion in Mercury's sky, when
Mercury is near perihelion.

Calculate the total duration of this apparent solar retrograde motion during
one orbit of Mercury around the Sun. Express your answer in Earth days.

\probhead{3.26}{Europa}{2023}{Katowice, Poland}{TH Q4}{10}
% source id: 2023_TH_Q4   tag: heat conduction through ice, tidal heating
\begin{description}
\item[(a)] Assuming that the ice covering the ocean on Europa is 6\,km thick,
  that the surface temperature on the night side of Europa is 100\,K and that
  the temperature at the ice-water boundary is 273\,K, calculate the total
  power corresponding to the heat emitted from the interior of this moon.
\item[(b)] On Earth, the mean geothermal heat flux measured at the
  continental surface is $70 \times 10^{-3}$\,W\,m$^{-2}$ and originates
  mainly from radioactive decay. Is the heat emanating from the interior of
  Europa more likely to come from radioactive decay or tidal forces? Assume
  that Earth and Europa have a similar isotopic composition. (Select the
  correct answer on the answer sheet and show your working.)
\end{description}

Notes: the heat passing through a wall with a surface area $S$ and thickness
$d$ in time $t$ is described by the formula:
\[ Q = \lambda S \Delta T t / d, \]
where $\lambda$ stands for thermal conductivity and $\Delta T$ for the
temperature difference.

The thermal conductivity of ice
$\lambda = 3$\,W\,m$^{-1}$K$^{-1}$. The mass and radius of Europa are
$4.8 \times 10^{22}$\,kg and 1561\,km.

\probhead{3.27}{Libration}{2023}{Katowice, Poland}{TH Q7}{20}
% source id: 2023_TH_Q7   tag: lunar libration, visible surface fraction
As a result of libration, studied among others by Johannes Hevelius, more
than half of the Moon's surface can be observed from Earth. Assume that the
observer is geocentric.

\begin{description}
\item[(a)] Estimate $\phi_B$, the maximum angle of libration in latitude. The
  axial tilt (obliquity) of the Moon with respect to its orbital plane is
  $\alpha = 6^\circ 41'$.
\item[(b)] Estimate $\phi_L$, the maximum angle of libration in longitude.
  Assume that the Moon is always aligned with the same side facing towards
  the second focus F2 of its orbit, and that the eccentricity of the Moon's
  orbit $e$ changes between 0.044 and 0.064 on a timescale of several months.
\item[(c)] Estimate the fraction of the Moon's surface which can be seen from
  Earth.
\item[(d)] Calculate how many months (lunations) are needed for an observer
  to see the Moon's surface determined in part (c).
\end{description}

\probhead{3.28}{Greatest Eclipse}{2024}{Rio de Janeiro, Brazil}{TH T10}{75}
% source id: 2024_TH_T10   tag: 1919 eclipse shadow geometry, totality
The greatest eclipse is defined as the instant when the axis of the Moon's
shadow cone gets closest to the centre of the Earth in a solar eclipse. This
problem explores the geometry of this phenomenon, using the solar eclipse of
29\textsuperscript{th} May 1919 as an example, as it has great historical
significance for being the first time astronomers were able to
observationally verify general relativity. One of the scientific expeditions
to observe this eclipse took place in the Brazilian city of Sobral.

The two following tables show the Cartesian and spherical coordinates of the
Sun and the Moon at the time of the greatest eclipse. The system used for
these coordinates is right-handed and has the origin at the centre of the
Earth, the positive $x$-axis pointing towards the Greenwich meridian, and the
positive $z$-axis pointing towards the North Pole. For the rest of this
problem, this will be referred to as \textbf{System I}.

Spherical Coordinates:
\begin{center}
\begin{tabular}{c|c|c}
 & Centre of the Sun & Centre of the Moon \\
\hline
Radial Distance ($r$) & $1.516 \times 10^{11}$\,m
  & $3.589 \times 10^{8}$\,m \\
Polar Angle ($\theta$) & $68^\circ 29' 44.1''$ & $68^\circ 47' 41.6''$ \\
Azimuthal Angle ($\varphi$) & $-1^{\mathrm{h}}11^{\mathrm{m}}28.2^{\mathrm{s}}$
  & $-1^{\mathrm{h}}11^{\mathrm{m}}22.9^{\mathrm{s}}$ \\
\end{tabular}
\end{center}

Cartesian Coordinates:
\begin{center}
\begin{tabular}{c|c|c}
 & Centre of the Sun & Centre of the Moon \\
\hline
$x$ & $1.342 \times 10^{11}$\,m & $3.185 \times 10^{8}$\,m \\
$y$ & $-4.327 \times 10^{10}$\,m & $-1.025 \times 10^{8}$\,m \\
$z$ & $5.557 \times 10^{10}$\,m & $1.298 \times 10^{8}$\,m \\
\end{tabular}
\end{center}

For this problem, assume that the Earth is a perfect sphere.

\textbf{Note:} The spherical coordinates of a point $P$ are defined as
follows:
\begin{itemize}
\item Radial distance ($r$): distance between the origin ($O$) and $P$
  (range: $r \geq 0$).
\item Polar angle ($\theta$): angle between the positive $z$-axis and the
  line segment $OP$ (range: $0^\circ \leq \theta \leq 180^\circ$)
\item Azimuthal angle ($\varphi$): angle between the positive $x$-axis and
  the projection of the line segment $OP$ onto the $xy$-plane (range:
  $-12^{\mathrm{h}} \leq \varphi < 12^{\mathrm{h}}$)
\end{itemize}

\ioaafig{2024_TH_T10-f1.pdf}
% point P, radius r, polar angle theta and azimuthal angle phi (page 10 of
% the 2024 theory paper, top of page)

\textbf{Part I: Geographic Coordinates (25 points)}

\begin{description}
\item[(a)] (3 points) Determine the declination of the Sun and the Moon
  during the greatest eclipse for a geocentric observer.
\item[(b)] (3 points) Determine the right ascension of the Sun and the Moon
  at the time of the greatest eclipse for a geocentric observer. The local
  sidereal time at Greenwich at that same moment was
  $5^{\mathrm{h}}32^{\mathrm{m}}35.5^{\mathrm{s}}$.
\item[(c)] (4 points) Find a unit vector that indicates the direction of the
  axis of the Moon's shadow cone. This vector should point from the Moon to
  the vicinity of the centre of the Earth.
\item[(d)] (15 points) Determine the latitude and the longitude of the point
  where the axis of the Moon's shadow cone crosses the surface of the Earth
  during the greatest eclipse.
\end{description}

\textbf{Part II: Duration of the Totality (50 points)}

Precisely determining the duration of totality of a solar eclipse involves
complex calculations that would be beyond the scope of this problem. However,
it is possible to obtain a reasonable approximation for this value using the
two following assumptions:
\begin{itemize}
\item The size of the umbra on the surface of the Earth remains roughly
  constant throughout totality at a given location.
\item The velocity of the umbra on the surface of the Earth remains roughly
  constant throughout totality at a given location.
\end{itemize}

\begin{description}
\item[(e)] (10 points) Estimate the radius of the umbra during the greatest
  eclipse. In order to simplify the calculations, assume that the umbra is
  small enough that it can be considered approximately flat and that the axis
  of the Moon's shadow cone is extremely close to the centre of the Earth
  during the greatest eclipse.
\item[(f)] (3 points) Calculate the velocity of the Earth's rotation at the
  latitude of the centre of the umbra.
\item[(g)] (4 points) Determine the orbital velocity of the Moon at the
  instant of the greatest eclipse. Neglect the changes in the semi-major axis
  of the Moon's orbit.
\end{description}

For the remaining items of this problem, assume that the tangential velocity
of the Moon is roughly the same as the orbital velocity and neglect its
radial component.

In order to calculate the velocity of the umbra, it is convenient to define
two new additional right-handed coordinate systems. \textbf{System II} is
defined as follows:
\begin{itemize}
\item Origin ($O_{\mathrm{II}}$): position of the Moon at the instant of the
  greatest eclipse.
\item Positive $x$-axis: Tangent to the declination circle. Points
  eastwards.
\item Positive $y$-axis: Tangent to the meridian of right ascension. Points
  northwards.
\end{itemize}

\textbf{System III} is defined as follows:
\begin{itemize}
\item Origin ($O_{\mathrm{III}}$): centre of the umbra at the instant of the
  greatest eclipse.
\item Positive $x$-axis: Tangent to the latitude circle. Points eastwards.
\item Positive $y$-axis: Tangent to the meridian of longitude. Points
  northwards.
\end{itemize}

Note that in both systems, the $xy$-plane is tangent to the celestial sphere
at the position of the origin.

System III is similar to System II, with the only difference being that the
origin ($O_{\mathrm{III}}$) is at the centre of the umbra at the moment of
the greatest eclipse.

\begin{description}
\item[(h)] (14 points) Using System II, determine the velocity vector of the
  Moon during the greatest eclipse. Note that the intersection between the
  Celestial Equator and the lunar orbit that is closer to the position of the
  eclipse has a right ascension of
  $23^{\mathrm{h}}07^{\mathrm{m}}59.2^{\mathrm{s}}$.
\item[(i)] (10 points) Write the velocity vector of the Moon in System III.
  Note that in System I, the azimuthal angle difference between the positions
  of the origins $O_{\mathrm{II}}$ and $O_{\mathrm{III}}$ is negligible, so
  you should only take into account the difference in the polar angles.
\item[(j)] (6 points) Calculate the speed of the centre of the umbra along
  the surface of the Earth at the instant of the greatest eclipse.
\item[(k)] (3 points) Estimate the duration of the totality of the eclipse at
  the location with the coordinates found in part (d).
\end{description}

\section*{Data-analysis problems}

\probhead{3.29}{Galilean moons}{2007}{Chiang Mai, Thailand}{DA D1}{4}
% source id: 2007_DA_D1   tag: identify Galilean moons from motion
Computer simulation of the planet Jupiter and its 4 Galilean moons is shown
on the screen similar to the view you may see through a small telescope.
After observing the movement of the moons, please identify the names of the
moons that appear at the end of the simulation. (Simulation will be played on
screen during the first fifteen minutes and the last fifteen minutes of the
exam)

\ioaafig{2007_DA_D1-f1.pdf}

\probhead{3.30}{Solar System objects}{2007}{Chiang Mai, Thailand}{DA D3}{8}
% source id: 2007_DA_D3   tag: identify objects from positions, elongations
A set of data containing the apparent positions of 4 Solar System objects
over a period of 1 calendar year is given in Table 1. Show your method of
data analysis carefully and answer the following questions.

\begin{center}
\begin{tabular}{lll}
Location of observer & Latitude: & N $18^\circ 47' 00.0''$ \\
 & Longitude: & E $98^\circ 59' 00.0''$ \\
\end{tabular}
\end{center}

\textit{[Table 1 not in source paper (distributed separately with exam)]}

\begin{description}
\item[3.1] Put the letters A, B, C and D beside the appropriate objects on
  the answer sheet. \textbf{(2 points)}
\item[3.2] During the period of observation, which object could be observed
  for the longest duration at night time? \textbf{(1 point)}
\item[3.3] What was the date corresponding to the situation in 3.2?
  \textbf{(1 point)}
\item[3.4] Assuming the orbits are coplanar (lie on the same plane) and
  circular, indicate the positions of the four objects and the Earth on the
  date in 3.3, in the orbit diagram provided in your answer sheet. The answer
  (sheet) must show one of the objects as the Sun at the centre of the Solar
  System. Other objects including the Earth must be specified together with
  the correct values of elongation on that date. \textbf{(4 points)}
\end{description}

\probhead{3.31}{The Age of Meteorite}{2008}{Bandung, Indonesia}{DA D3}{100}
% source id: 2008_DA_D3   tag: Rb-Sr isochron regression
The basic equation of radioactive decay can be expressed as:
\[ N(t) = N_0 \exp(-\lambda t) \]
where $N(t)$ and $N_0$ are the number of remaining atoms of the radioactive
isotope (or parent isotope) at time $t$ and its initial number at $t = 0$,
respectively, while $\lambda$ is the decay constant. The decay of the parent
produces daughter nuclides $D(t)$, or radiogenics, which is defined as
\[ D(t) = N_0 - N(t)\,. \]

Based on those ideas, a group of astronomers investigates a number of
meteorite samples to determine their ages. They have two kinds of samples:
allende chondrite (A) and basaltic achondrite (B). From the samples, they
measure the abundance of $^{87}$Rb and $^{87}$Sr, where it is assumed that
$^{87}$Sr is entirely produced by the decay of $^{87}$Rb. The value of
$\lambda$ is $1.42 \times 10^{-11}$ per year for this isotopic decay. In
addition, non-radiogenic element $^{86}$Sr is also measured. Results of
measurement are given in the table below, expressed in ppm (part per
million).

\begin{center}
\begin{tabular}{c|c|c|c|c}
Sample No & Meteorite type & $^{86}$Sr (ppm) & $^{87}$Rb (ppm)
  & $^{87}$Sr (ppm) \\
\hline
1  & A & 29.6 & 0.3   & 20.7 \\
2  & B & 58.7 & 68.5  & 44.7 \\
3  & B & 74.2 & 14.4  & 52.9 \\
4  & A & 40.2 & 7.0   & 28.6 \\
5  & A & 19.7 & 0.4   & 13.8 \\
6  & B & 37.9 & 31.6  & 28.4 \\
7  & A & 33.4 & 4.0   & 23.6 \\
8  & B & 29.8 & 105.0 & 26.4 \\
9  & A & 9.8  & 0.8   & 6.9  \\
10 & B & 18.5 & 44.0  & 15.4 \\
\end{tabular}
\end{center}

Questions:
\begin{enumerate}
\item Express the time $t$ in term of $\dfrac{D(t)}{N(t)}$.
\item Determine the half-life $t_{1/2}$, i.e., the time needed to obtain a
  half number of parents after decay.
\item Knowledge on the ratio between two isotopes is more valuable than just
  the absolute abundance of each isotope. It is quite likely that there was
  some initial strontium present. By taking
  $\left(\frac{^{87}\mathrm{Rb}}{^{86}\mathrm{Sr}}\right)$ as independent
  variable and $\left(\frac{^{87}\mathrm{Sr}}{^{86}\mathrm{Sr}}\right)$ as
  dependent variable, estimate the simple linear regression model to
  represent the data.
\item Plot $\left(\frac{^{87}\mathrm{Rb}}{^{86}\mathrm{Sr}}\right)$ versus
  $\left(\frac{^{87}\mathrm{Sr}}{^{86}\mathrm{Sr}}\right)$ and also the
  regression line (isochrone) for each type of the meteorites. (Please use
  minimum 7 decimal digits for intermediate calculations)
\item Subsequently, help this astronomer to determine the age of each type of
  the meteorites and its error. Which type is older?
\item Determine the initial value of
  $\left(\frac{^{87}\mathrm{Sr}}{^{86}\mathrm{Sr}}\right)_0$ for each type of
  the meteorites and its error.
\end{enumerate}

\textit{Glossary}:

A simple linear regression line $y = a + bx$ can be fitted to a set of data
$(X_i, Y_i)$, $i = 1, \ldots, n$, in which
\[ b = \frac{SS_{xy}}{SS_{xx}} \]
\[ a = \bar{y} - b\bar{x} \]
where
\[ SS_{xx} \text{ : sum of square for } X
   = \sum_{i=1}^{n} X_i^2
   - \frac{1}{n}\left( \sum_{i=1}^{n} X_i \right)^2 \]
\[ SS_{yy} \text{ : sum of square for } Y
   = \sum_{i=1}^{n} Y_i^2
   - \frac{1}{n}\left( \sum_{i=1}^{n} Y_i \right)^2 \]
\[ SS_{xy} \text{ : sum of square for both } X \text{ and } Y
   = \sum_{i=1}^{n} X_i Y_i
   - \frac{1}{n} \sum_{i=1}^{n} X_i \sum_{i=1}^{n} Y_i \]

Standard deviation of each parameter, $a$ and $b$ can be calculated by
\[ S_a = \sqrt{ \frac{SS_{YY} - \dfrac{(SS_{XY})^2}{SS_{XX}}}
   {(n-2)\,SS_{XX}} \times \sum_{i=1}^{n} X_i^2 } \]
\[ S_b = \sqrt{ \frac{SS_{YY} - \dfrac{(SS_{XY})^2}{SS_{XX}}}
   {(n-2)\,SS_{XX}} } \]

\probhead{3.32}{Venus}{2009}{Tehran, Iran}{DA D2}{60}
% source id: 2009_DA_D2   tag: Venus orbit elements from phase
An observer in Deh-Namak (you will be taking the observational part of the
exam in this region tonight) has observed Venus for seven months, started
from September 2008 and continued until March 2009. During the observation, a
research CCD camera and an image processing software were used to take high
resolution images and to extract high precision data. Table 2.1 shows the
collected data during the observation.

Table 2.1 description:
\begin{center}
\begin{tabular}{lp{9cm}}
Column 1 & Date of observation. \\
Column 2 & Earth--Sun distance in astronomical unit (AU) for observation
  date and time. This value is taken from high precision tables. \\
Column 3 & Phase of Venus, Percent of Venus disk illuminated by the Sun as
  observed from the Earth. \\
Column 4 & Elongation of Venus, the angular distance between center of the
  Sun and center of Venus in degrees as observed from the Earth. \\
\end{tabular}
\end{center}

\begin{center}
Table 2.1

\begin{tabular}{c|c|c|c}
Column 1 & Column 2 & Column 3 & Column 4 \\
\hline
Date & Earth--Sun Distance (AU) & Phase (\%)
  & Elongation (SEV) ($^\circ$) \\
\hline
20/9/2008   & 1.0043 & 88.4 & 27.56 \\
10/10/2008  & 0.9986 & 84.0 & 32.29 \\
20/10/2008  & 0.9957 & 81.6 & 34.53 \\
30/10/2008  & 0.9931 & 79.0 & 36.69 \\
9/11/2008   & 0.9905 & 76.3 & 38.71 \\
19/11/2008  & 0.9883 & 73.4 & 40.62 \\
29/11/2008  & 0.9864 & 70.2 & 42.38 \\
19/12/2008  & 0.9839 & 63.1 & 45.29 \\
29/12/2008  & 0.9834 & 59.0 & 46.32 \\
18/1/2009   & 0.9838 & 49.5 & 47.09 \\
7/2/2009    & 0.9863 & 37.2 & 44.79 \\
17/2/2009   & 0.9881 & 29.6 & 41.59 \\
27/2/2009   & 0.9904 & 20.9 & 36.16 \\
19/3/2009   & 0.9956 & 3.8  & 16.08 \\
\end{tabular}
\end{center}

\begin{description}
\item[a)] Using given data in table 2.1, calculate the Sun-Venus-Earth angle
  ($\angle SVE$). This is angular separation between the Sun and the Earth as
  seen from Venus. Write $\angle SVE$ angle in column 2 of Table 2.2 in your
  answer sheet for the all observing dates.

  Hint: Remember that the line between light and shadow, in the phases, is an
  ellipse arc.

\ioaafig{2009_DA_D2-f1.pdf}
\begin{center}
Figure (part a)
\end{center}

\item[b)] Calculate Sun--Venus distance in AU and write it down in column 3
  of table 2.2 for all observation dates.
\item[c)] Plot Sun--Venus distance versus observing date.
\item[d)] Find perihelion ($r_{v,min}$) and aphelion ($r_{v,max}$) distances
  of Venus from the Sun.
\item[e)] Calculate semi-major axis ($a$) of the Venus orbit.
\item[f)] Calculate eccentricity ($e$) of Venus orbit.
\end{description}

\probhead{3.33}{Asteroid photometry}{2012}{Rio de Janeiro, Brazil}{DA D1}{}
% source id: 2012_DA_D1   tag: asteroid light curves, phase curve, regolith
A few facts about photometry of asteroids:
\begin{itemize}
\item Asteroids are small, irregularly shaped objects of our Solar System
  that orbit the Sun in approximately elliptical orbits.
\item Their brightness observed at a given instant from Earth depend on the
  surface area illuminated by the Sun and the part of the asteroid which is
  visible to the observer. Both vary as the asteroid moves.
\item The way the sunlight is reflected by the surface of the asteroid
  depends on its texture and on the angle between the Sun, the asteroid and
  the observer (phase angle), which varies as the Earth and the asteroid move
  along their orbits. In particular, asteroids with surfaces covered by fine
  dust (called regoliths) exhibit a sharp increase in brightness at phase
  angles $\varphi$ close to zero (i.e., when they are close to the
  opposition).
\item Since the observed flux of any source decreases with the square of
  distance, the observed magnitude of an asteroid also depends on its
  distance from the Sun and from the observer at the time of the observation.
  Their apparent magnitude $m$ outside the atmosphere is then
  \[ M(t) = M_r(t) + 5\log(RD) \]
  where $m_r$ is usually called reduced magnitude (meaning the magnitude the
  asteroid would have if its distances from the Sun and the Earth were
  reduced both to 1\,AU) and depends only on the visible illuminated surface
  area and on phase angle effects. $R$ and $D$ are, respectivelly, the
  % sic: "respectivelly" as in the original
  heliocentric and geocentric distances.
\end{itemize}

Consider now the following scenario. Light curves of a given asteroid were
obtained at three different nights at different points of its orbit, and at
each time a photometric standard was observed in the same frame of the
asteroid. Table 1 shows the geometric configuration of the asteroid at each
night (phase angle $\varphi$, in degrees, $R$ and $D$ in AUs), and the
calibrated magnitude of the photometric standard star that was observed along
with the asteroid. Consider the calibrated magnitude as the final apparent
magnitude after all the effects are taken into account.

Tables 2, 3 and 4 contain, for each night, the time interval of each
observation with respect to the first one (in hours) the air mass, the
instrumental magnitude of the asteroid and the instrumental magnitude of the
star.

Air mass is the dimensionless thickness of the atmosphere along the line of
sight, and is equal to 1 when looking to zenith.

\begin{center}
Table 1

\begin{tabular}{c|c|c|c|c}
Night & $D$ & $R$ & $\varphi$ & $M_{\mathrm{star}}$ \\
\hline
A & 0.36 & 1.35 & 0.0  & 8.2 \\
B & 1.15 & 2.13 & 8.6  & 8.0 \\
C & 2.70 & 1.89 & 15.6 & 8.1 \\
\end{tabular}
\end{center}

\begin{center}
Table 2: Night A

\begin{tabular}{c|c|c|c}
$\Delta t$ & Air mass & $m_{\mathrm{ast}}$ & $m_{\mathrm{star}}$ \\
\hline
0.00 & 1.28 & 7.44 & 8.67 \\
0.44 & 1.18 & 7.38 & 8.62 \\
0.89 & 1.11 & 7.34 & 8.59 \\
1.33 & 1.06 & 7.28 & 8.58 \\
1.77 & 1.02 & 7.32 & 8.58 \\
2.21 & 1.00 & 7.33 & 8.56 \\
2.66 & 1.00 & 7.33 & 8.56 \\
3.10 & 1.01 & 7.30 & 8.56 \\
3.54 & 1.03 & 7.27 & 8.58 \\
3.99 & 1.07 & 7.27 & 8.58 \\
4.43 & 1.13 & 7.31 & 8.61 \\
4.87 & 1.21 & 7.37 & 8.63 \\
5.31 & 1.32 & 7.42 & 8.67 \\
5.76 & 1.48 & 7.49 & 8.73 \\
6.20 & 1.71 & 7.59 & 8.81 \\
6.64 & 2.06 & 7.69 & 8.92 \\
7.09 & 2.62 & 7.87 & 9.14 \\
7.53 & 3.67 & 8.21 & 9.49 \\
\end{tabular}
\end{center}

\begin{center}
Table 3: Night B

\begin{tabular}{c|c|c|c}
$\Delta t$ & Air mass & $m_{\mathrm{ast}}$ & $m_{\mathrm{star}}$ \\
\hline
0.00 & 1.28 & 13.24 & 8.38 \\
0.44 & 1.18 & 13.21 & 8.36 \\
0.89 & 1.11 & 13.13 & 8.34 \\
1.33 & 1.06 & 13.11 & 8.33 \\
1.77 & 1.02 & 13.11 & 8.32 \\
2.21 & 1.00 & 13.15 & 8.32 \\
2.66 & 1.00 & 13.17 & 8.32 \\
3.10 & 1.01 & 13.17 & 8.32 \\
3.54 & 1.03 & 13.13 & 8.33 \\
3.99 & 1.07 & 13.15 & 8.34 \\
4.43 & 1.13 & 13.14 & 8.34 \\
4.87 & 1.21 & 13.14 & 8.37 \\
5.31 & 1.32 & 13.21 & 8.38 \\
5.76 & 1.48 & 13.30 & 8.43 \\
6.20 & 1.71 & 13.34 & 8.47 \\
6.64 & 2.06 & 13.39 & 8.54 \\
7.09 & 2.62 & 13.44 & 8.65 \\
7.53 & 3.67 & 13.67 & 8.87 \\
\end{tabular}
\end{center}

\begin{center}
Table 4: Night C

\begin{tabular}{c|c|c|c}
$\Delta t$ & Air mass & $m_{\mathrm{ast}}$ & $m_{\mathrm{star}}$ \\
\hline
0.00 & 1.28 & 11.64 & 8.58 \\
0.44 & 1.18 & 11.53 & 8.54 \\
0.89 & 1.11 & 11.56 & 8.60 \\
1.33 & 1.06 & 11.49 & 8.52 \\
1.77 & 1.02 & 11.58 & 8.48 \\
2.21 & 1.00 & 11.79 & 8.63 \\
2.66 & 1.00 & 11.67 & 8.53 \\
3.10 & 1.01 & 11.53 & 8.46 \\
3.54 & 1.03 & 11.47 & 8.48 \\
3.99 & 1.07 & 11.63 & 8.67 \\
4.43 & 1.13 & 11.51 & 8.51 \\
4.87 & 1.21 & 11.65 & 8.55 \\
5.31 & 1.32 & 11.77 & 8.61 \\
5.76 & 1.48 & 11.88 & 8.75 \\
6.20 & 1.71 & 11.86 & 8.78 \\
6.64 & 2.06 & 12.03 & 9.03 \\
7.09 & 2.62 & 12.14 & 9.19 \\
7.53 & 3.67 & 12.63 & 9.65 \\
\end{tabular}
\end{center}

\begin{enumerate}
\item Plot the star magnitude versus air mass for each set
\item Calculate the extinction coefficient for each night (see Appendix A at
  the end of the text)
\item Were the observations affected by clouds in one night? Express your
  answer as: a) Night A, b) Night B, c) Night C, d) None of the nights.
\item Plot the calibrated magnitude versus time for each set of observations
  of the asteroid (see Appendix B).
\item Determine the rotation period for each night. Consider that the light
  curve for this asteroid has two minima and two maxima, and that the
  semi-period is the average of the intervals between the two maxima and the
  two minima.
\item Determine the amplitude (difference from maximum to minimum) of the
  light curve for each night
\item Plot the calibrated reduced magnitude $M_r$ versus phase angle
  $\varphi$ (use the mean value of each light curve)
\item Calculate the angular coefficient of the phase curve (the plot of the
  calibrated reduced magnitude versus the phase angle) considering only the
  points away from the opposition (See 3rd bullet of facts about photometry
  of asteroids above)
\item Is there any reason to assume a surface covered by fine dust (regolith)?
  Answer YES/NO.
\end{enumerate}

\textbf{Appendix A: extinction coefficients}

The magnitude outside the atmosphere is given by:
\[ M = m - A.X + B \]
where $A$ is the extinction coefficient and $B$ is the zero point coefficient
for the night, $X$ is the airmass of the observation and $m$ is the
instrumental magnitude (that is, the magnitude obtained directly from the
image)

\textbf{Appendix B: calibrated differential magnitude}

If the standard star is the same frame as the object, the calibrated
magnitude of the object can be obtained by:
\[ M_{ast} = m_{ast} - m_{star} + M_{star} \]

\textbf{Appendix C: least squares estimation of angular coefficients}

Given straight lines described by
\[ y_i = \alpha + \beta x_i, \quad i = 1..n \]
the least squares estimator of $\beta$ is given by
\[ \beta = \frac{\sum_i^n x_i y_i
     - \frac{1}{n} \sum_i^n x_i \sum_i^n y_i}
   {\sum_i^n x_i^2 - \frac{1}{n} \left( \sum_i^n x_i \right)^2} \]

\probhead{3.34}{Thermodynamic test}{2014}{Suceava, Romania}{DA D2}{}
% source id: 2014_DA_D2   tag: CO2 atmosphere profiles from probe
An hypotetical shuttle is launched to investigate the atmosphere (100\%
% sic: "An hypotetical" as in the original
CO$_2$) of two extrasolar planets P$_1$ and P$_2$. The atmosphere is in
static thermodynamic equilibrium. When the shuttle is near each planet, a
radio probe is launched toward respective planet, in vertical direction (in
the direction of the planet's radius). When the radio probe reaches constant
velocity, it starts sending values of the pressure of the atmosphere. In
Fig.~3.1 is plotted the atmospheric pressure values (in arbitrary units) as
function of the time of descent for the planet P$_1$. When the probe touches
the surface of planet P$_1$ it sends the value of the temperature
$T_0 = 700$\,K and the value of the gravitational acceleration
$g_0 = 10\,\mathrm{m\,s}^{-2}$.

\ioaafig{2014_DA_D2-s7.pdf}
\begin{center}
Fig.~3.1.
\end{center}

The gravitational acceleration on each planet is assumed to be constant
during uniform descent of the radio probes.

\begin{description}
\item[a)] Find the altitude $h_0$ from where the radio probe R$_1$ starts the
  uniform descent and thus starts the transmitting information.
\item[b)] Find the temperature of planet P$_1$ at the altitude
  $h = 39.6$\,km. You know: The universal gas constant
  $R = 8.3$\,J/(mol\,K); the molar mass of CO$_2$, $\mu = 44$\,g/mol.
\item[c)] In Fig.~3.2.\ was plotted the atmospheric pressure values (in
  arbitrary units) as a function of time of descent for the planet P$_2$
  atmosphere. When the probe touched the surface of the planet P$_2$, it
  sends the value of the temperature $T_0 = 750$\,K and respectively the
  value of gravitational acceleration $g = 8\,\mathrm{m\,s}^{-2}$.

  Draw the following dependency graphs for $p = f(h)$ and $T = f(h)$ in the
  CO$_2$ atmosphere of the planet P$_2$.
\end{description}

\ioaafig{2014_DA_D2-s8.pdf}
\begin{center}
Fig.~3.2.
\end{center}

\probhead{3.35}{Distance to the Moon}{2016}{Bhubaneswar, India}{DA D2}{50}
% source id: 2016_DA_D2   tag: lunar ephemerides, eclipse umbra
Geocentric ephemerides of the Moon for September 2015 are given in the form
of a table. Each reading was taken at 00:00 UT.

\begin{center}
\begin{tabular}{c|ccc|ccc|c|c|c}
Date & \multicolumn{3}{c|}{R.A. ($\alpha$)}
  & \multicolumn{3}{c|}{Dec. ($\delta$)}
  & Angular Size ($\theta$) & Phase ($\phi$) & Elongation \\
 & h & m & s & $^\circ$ & $'$ & $''$ & $''$ & & Of Moon \\
\hline
Sep 01 & 0 & 36 & 46.02 & 3  & 6  & 16.8 & 1991.2 & 0.927 & $148.6^\circ$ W \\
Sep 02 & 1 & 33 & 51.34 & 7  & 32 & 26.1 & 1974.0 & 0.852 & $134.7^\circ$ W \\
Sep 03 & 2 & 30 & 45.03 & 11 & 25 & 31.1 & 1950.7 & 0.759 & $121.1^\circ$ W \\
Sep 04 & 3 & 27 & 28.48 & 14 & 32 & 4.3  & 1923.9 & 0.655 & $107.9^\circ$ W \\
Sep 05 & 4 & 23 & 52.28 & 16 & 43 & 18.2 & 1896.3 & 0.546 & $95.2^\circ$ W \\
\end{tabular}
\end{center}
% TABLE ABRIDGED — full data (Sep 01 to Sep 30) in crop edition; the complete
% table is also wired below as a figure.
\ioaafig{2016_DA_D2-f1.pdf}

The composite graphic\footnote{Credit: NASA's Scientific Visualization
Studio} below shows multiple snapshots of the Moon taken at different times
during the total lunar eclipse, which occurred in this month. For each shot,
the centre of frame was coinciding with the central north-south line of
umbra.

For this problem, assume that the observer is at the centre of the Earth and
angular size refers to angular diameter of the object / shadow.

\ioaafig{2016_DA_D2-f2.pdf}

\begin{description}
\item[(D2.1)] In September 2015, apogee of the lunar orbit is closest to

  New Moon / First Quarter / Full Moon / Third Quarter.

  Tick the correct answer in the Summary Answersheet. No justification for
  your answer is necessary. \hfill[3]
\item[(D2.2)] In September 2015, the ascending node of lunar orbit with
  respect to the ecliptic is closest to

  New Moon / First Quarter / Full Moon / Third Quarter.

  Tick the correct answer in the Summary Answersheet. No justification for
  your answer is necessary. \hfill[4]
\item[(D2.3)] Estimate the eccentricity, $e$, of the lunar orbit from the
  given data. \hfill[4]
\item[(D2.4)] Estimate the angular size of the umbra,
  $\theta_{\mathrm{umbra}}$, in terms of the angular size of the Moon,
  $\theta_{\mathrm{Moon}}$. Show your working on the image given on the
  backside of the Summary Answersheet. \hfill[8]
\item[(D2.5)] The angle subtended by the Sun at Earth on the day of the lunar
  eclipse is known to be $\theta_{\mathrm{Sun}} = 1915.0''$. In the figure
  below, $S_1 R_1$ and $S_2 R_2$ are rays coming from diametrically opposite
  ends of the solar disk. The figure is not to scale.

\ioaafig{2016_DA_D2-f3.pdf}

  Calculate the angular size of the penumbra, $\theta_{\mathrm{penumbra}}$,
  in terms of $\theta_{\mathrm{Moon}}$. Assume the observer to be at the
  centre of the Earth. \hfill[9]
\item[(D2.6)] Let $\theta_{\mathrm{Earth}}$ be angular size of the Earth as
  seen from the centre of the Moon. Calculate the angular size of the Moon,
  $\theta_{\mathrm{Moon}}$, as would be seen from the centre of the Earth on
  the eclipse day in terms of $\theta_{\mathrm{Earth}}$. \hfill[5]
\item[(D2.7)] Estimate the radius of the Moon, $R_{\mathrm{Moon}}$, in km
  from the results above. \hfill[3]
\item[(D2.8)] Estimate the shortest distance, $r_{\mathrm{perigee}}$, and the
  farthest distance, $r_{\mathrm{apogee}}$, to the Moon. \hfill[4]
\item[(D2.9)] Use appropriate data from September 10 to estimate the
  distance, $d_{\mathrm{Sun}}$, to the Sun from the Earth. \hfill[10]
\end{description}

\section*{Observation --- planetarium}

\probhead{3.36}{Hypothesized Planet in the Solar System}{2025}{Mumbai, India}{OBS-PLAN OP01}{25}
% source id: 2025_OBS-PLAN_OP01   tag: orbit inclinations from planetarium traces
For this question use sky map Map-OP01. This is an angle preserving
projection of the sky.

The sky projected on the planetarium dome is as seen from the solar system
barycentre. The Earth's orbit as seen from this position has been marked in
green both on the projected sky and in sky map Map-OP01. The first 3 minutes
are for dark adaptation and familiarization of the sky.

After first 3 minutes, four Trans-Neptunian Objects (TNOs) labelled O1, O2,
O3 and O4 will begin to appear on the dome. The brightness of the TNOs has
been enhanced to make them visible. No other solar system object is visible
in the projected sky. Time has been sped up to make their movements
noticeable.

We will assume that the orbital plane of a hypothesized new planet has an
inclination equal to the mean of the inclinations of the orbits of these 4
TNOs.

\ioaafig{2025_OBS-PLAN_OP01-s3.pdf}

\begin{description}
\item[(OP01.1)] Arrange the 4 TNOs in ascending order of the lengths of their
  semi-major axes and write their names in the appropriate boxes in the
  Summary Answersheet. Show measurements in the working sheet to justify your
  answer. \hfill[8]
\item[(OP01.2)] Trace at least 50\% of the visible trajectories for each of
  the 4 TNOs on the map Map-OP01. From these traces find the inclination
  angles, $i_1, i_2, i_3, i_4$, of the 4 TNOs, respectively, with respect to
  the Ecliptic. Hence obtain the best estimate of the inclination,
  $i_{\mathrm{hp}}$, of the plane of the orbit of the hypothesized planet.
  \hfill[17]
\end{description}

\probhead{3.37}{Observer on a Ringed Planet}{2025}{Mumbai, India}{OBS-PLAN OP02}{25}
% source id: 2025_OBS-PLAN_OP02   tag: sky as seen from a ringed planet
For this question use sky map Map-OP02.

The sky projected on the planetarium dome is as seen by an observer on some
hypothetical ringed planet. The observer is located very near to the equator
of the planet.

Assume that the ring lies in the equatorial plane of the planet and the
distance of the planet from its parent star is more than 15\,au. Note that
except for the 10 moons of this planet no other object from this planetary
system (other planets, their satellites, asteroids and even the parent star)
are visible on the projected sky.

At the start of the question the sky will make 3 full rotations at the rate
of 1 rotation every minute. After that the sky will stop rotating at the
observer's midnight. This sky view will remain on the dome for next 7
minutes.

Solve the following questions once the sky stops rotating.

\ioaafig{2025_OBS-PLAN_OP02-f1.pdf}
% exam; not part of the question paper PDF.

\begin{description}
\item[(OP02.1)] Mark the positions of any 4 (only 4) out of the 10 visible
  moons of the planet with a $\otimes$ sign and label them as ``M'' in the
  sky map. \hfill[8]
\item[(OP02.2)] The figure shown below is as seen from the top of the north
  pole of the planet (grey circle) and the horizontal lines are parallel
  light rays from its parent star.

% FIGURE NEEDED: diagram from page 2 of the 2025 planetarium paper — grey
% circle of radius R_P = 6.00 x 10^4 km labelled P, with parallel light rays
% from the parent star drawn as horizontal lines.

  Using the view of the ring on the dome, draw the inner and outer boundaries
  of the ring on this diagram, reproduced in the Summary Answersheet. Your
  sketch must be to scale. \hfill[12]
\item[(OP02.3)] Use the diagram above to determine the width of the ring,
  $w_{\mathrm{ring}}$ (in km), given that the radius of the planet is
  $6.00 \times 10^4$\,km. \hfill[5]
\end{description}

\clearpage

\probhead{3.38}{Retrograde Mars}{2023}{Katowice, Poland}{OBS-PLAN 2}{20}
% source id: 2023_OBS-PLAN_2   tag: Mars retrograde loop, heliocentric geometry
The projector will display Mars moving relative to the background stars over
one season of visibility (1.5 years) starting from the heliacal rising, chosen
so that Mars will be at maximum ecliptic latitude at opposition.

The ecliptic will also be displayed, marked with the positions of the Sun
during the year and the current date. The Sun will always be below the
horizon. Synodic period of Mars $= 780$ days. (Projection time: 10 minutes.)

\begin{description}
\item[(a)] Record the following quantities: \hfill(8 points)
  \begin{description}
  \item[i.] the dates of quadrature (when the elongation of Mars is $90^\circ$)
  \item[ii.] the date of the beginning of retrograde motion and the date of
    the end of retrograde motion
  \item[iii.] the date of opposition
  \item[iv.] the ecliptic latitude at opposition
  \item[v.] the width in ecliptic longitude of the loop made by the planet
  \end{description}
\end{description}

Based on your observations and assuming the orbits of Earth and Mars are
circular,
\begin{description}
\item[(b)] On the answer sheet, mark the positions of the Sun, Earth and Mars
  at the moments of opposition and quadrature in the heliocentric system and
  determine the radius of the orbit of Mars in a.u.\ geometrically, without
  using Kepler's Laws. Show your method in the answer sheet. \hfill(9 points)
\item[(c)] Derive the inclination of the orbit of Mars to the ecliptic.
  \hfill(3 points)
\end{description}

\section*{Solutions to Chapter 3}

\solhead{3.1}{}
% source id: 2007_TH_1.2
\ioaafig{2007_TH_1.2-s1.pdf}
The angular separation is maximum when Sun, Venus and Earth form a
right-angled triangle as shown. Here
\begin{align*}
R &= R_\oplus \sin 46^\circ\\
  &= (1\,\mathrm{A.U.})\sin 46^\circ\\
  &= 0.72\ \mathrm{A.U.}
\end{align*}

\solhead{3.2}{}
% source id: 2007_TH_1.4
During a full Moon we see the \underline{whole face} of the Moon. Hence
\[
\frac{\text{Moon's diameter}}{\text{distance to Moon}}
  = \text{angle in radians}
  = \frac{0.46 \times \pi}{180}
\]
\begin{align*}
\text{Distance to the Moon}
  &= \left(\text{Moon's diameter}\right)\times\frac{180}{0.46\times\pi}\\
  &= \left(2\times 1.7374\times 10^{6}\ \mathrm{m}\right)
     \times\frac{180}{0.46\times\pi}\\
  &= 4.328\times 10^{8}\ \mathrm{m} \;=\; 4.3\times 10^{5}\ \mathrm{km}
\end{align*}

\solhead{3.3}{A Planet \& Its Surface Temperature}
% source id: 2007_TH_Q2
\begin{description}
\item[a.)] Intensity
  \[ I = \frac{L}{4\pi D^2} \]
  \hfill(1 point)
\item[b.)] Absorption rate
  \begin{align*}
  \mathcal{A} &= \left(1-\alpha\right)\pi R^2 I\\
              &= \left(1-\alpha\right)\frac{LR^2}{4D^2}
  \end{align*}
  \hfill(1 point)
\item[c.)] Light energy reflected by the planet per unit time is
  \[ \alpha\pi R^2 I = \frac{\alpha L R^2}{4D^2} \]
  \hfill(1 point)

  Hence the planet's luminosity is
  \[ L_{\mathrm{planet}} = \frac{\alpha L R^2}{4D^2} \]
  \hfill(1 point)
\item[d.)] Here, we will neglect the planet's internal source of energy. Let
  $T$ be the black-body temperature of the planet's surface in kelvins. Since
  the planet is rotating fast, we may assume that its surface is being heated
  up uniformly to approximately the same temperature $T$. The total amount of
  black-body radiation emitted by the planet's surface is from
  Stefan-Boltzmann law: $4\pi R^2 \cdot \sigma T^4$, $\sigma$ being
  Stefan-Boltzmann constant. At equilibrium, that is when the temperature
  remains steady, this emission rate must be equal to the absorption rate in
  b.). \hfill(1 point)

  Hence
  \begin{align*}
  4\pi R^2 \cdot \sigma T^4 &= \left(1-\alpha\right)\frac{LR^2}{4D^2}\\
  T &= \left[\left(1-\alpha\right)\frac{L}{16\pi\sigma D^2}\right]^{\frac14}
  \end{align*}
  \hfill(1 point)
\item[e.)] In this case the emitted black-body radiation is mostly from the
  planet's surface facing the star. The emitting surface area is now only
  $2\pi R^2$ and not $4\pi R^2$. Hence the surface temperature is given by
  $T'$, where
  \[ 2\pi R^2 \cdot \sigma \left(T'\right)^4
     = \left(1-\alpha\right)\frac{LR^2}{4D^2} \]
  \hfill(1 point)
  \[ T' = \left[\left(1-\alpha\right)\frac{L}{8\pi\sigma D^2}\right]^{\frac14}
        = \left(2\right)^{\frac14}\cdot T \approx \left(1.19\right)T \]
  \hfill(1 point)
\item[f.)]
  \begin{align*}
  T &= \left[\left(1-\alpha\right)\frac{L}{16\pi\sigma D^2}\right]^{\frac14}\\
  T &= \left[\left(1-0.25\right)\times
       \frac{3.826\times 10^{26}}
            {16\pi\times 5.67\times 10^{-8}\times
             \left(1.523\times 1.496\times 10^{11}\right)^2}\right]^{\frac14}\\
    &= 209.8 \simeq 210\ \mathrm{K} = -63\,^\circ\mathrm{C}
  \end{align*}
  \hfill(2 points)
\end{description}

\solhead{3.4}{}
% source id: 2008_TH_S11
$D = 10$\,cm, $f$-ratio 15 or $f = F/D = 15$, then $F = 150$\,cm
($D$ = diameter of the telescope and $F$ is the focal length of the
telescope).

$S$ = scale of the image on the focal plane of the telescope
$= (206265/F)\ (''/\mathrm{mm})$
\[
= (206265/1500) = 137.51\ ''/\mathrm{mm}
\]
\hfill(25\%)

Angular diameter of the Sun
\begin{align*}
&= 137.51\ ''/\mathrm{mm} \times 13.817\ \mathrm{mm}\\
&= 1899.97567'' = 0^\circ.527771019 = 0^\circ\,31'\,39''.98
\end{align*}
Distance of the Sun
\begin{align*}
&= 1\,392\,000/\left((0^\circ.527771019/180^\circ)\times\pi\right)\\
&= 151\,118\,045.9\ \mathrm{km}
\end{align*}
\hfill(25\%)

Angular diameter of the Moon
\begin{align*}
&= 137.51\ ''/\mathrm{mm} \times 13.235\ \mathrm{mm}\\
&= 1819.94485'' = 0^\circ.505540236 = 0^\circ\,30'\,19''.94
\end{align*}
Distance of the Moon
$= 3476/(0^\circ.505540236/180^\circ \times \pi) = 393955.0515$\,km.
\hfill(25\%)

The percentage of the solar disk covered by the Moon is
$(13.235\ \mathrm{mm}/13.817\ \mathrm{mm})^2 \times 100\% = 91.75\,\%$.
\hfill(25\%)

\solhead{3.5}{}
% source id: 2008_TH_S13
\textbf{I. Algebraic Solution}

\ioaafig{2008_TH_S13-s1.pdf}
Let $y = mx$ be the line of sight of the telescope. \hfill(20\%)

Then, the intersection of the line of sight and the circle is
\begin{align*}
x^2 + (mx)^2 - 10x - 8(mx) + 40 &= 0\\
\left(1+m^2\right)x^2 - (10+8m)x + 40 &= 0
\end{align*}
\hfill(20\%)

The equation will have solution if its discriminant
\begin{align*}
D &= b^2 - 4ac \ge 0,\\
  &= (-(10+8m))^2 - 4(1+m^2)40 \ge 0.
\end{align*}
The solution of the inequality is
\[
\frac{5}{6} - \frac{\sqrt{10}}{12} \le m \le
\frac{5}{6} + \frac{\sqrt{10}}{12}
\]
\hfill(20\%)

Therefore the exact value of the maximum value of $\tan\varphi$ is
$\dfrac{5}{6} + \dfrac{\sqrt{10}}{12}$ \hfill(20\%)

and the minimum value is $\dfrac{5}{6} - \dfrac{\sqrt{10}}{12}$ \hfill(20\%)

\textbf{II. Geometric Solution}

\ioaafig{2008_TH_S13-s2.pdf}
Let us write the equation of the circle in the following form \hfill(20\%)
\begin{align*}
(x-5)^2 + (y-4)^2 &= -40 + 25 + 16\\
                  &= 1
\end{align*}
The centre of the circle is $(5,4)$ and the radius is 1.

First we obtain $\tan\alpha = \dfrac{4}{5}$ and
$\tan\beta = \dfrac{1}{\sqrt{40}}$ \hfill(20\%)

The minimum value of tan of elevation angle is \hfill(30\%)
\[
\tan(\alpha-\beta) = \frac{\tan\alpha - \tan\beta}{1 + \tan\alpha\tan\beta}
 = \frac{\dfrac{4}{5} - \dfrac{1}{\sqrt{40}}}{1 + \dfrac{4}{5\sqrt{40}}}
 = \frac{5}{6} - \frac{\sqrt{10}}{12}
\]
The maximum value of tan of elevation angle is
\[
\tan(\alpha+\beta) = \frac{\tan\alpha + \tan\beta}{1 - \tan\alpha\tan\beta}
 = \frac{\dfrac{4}{5} + \dfrac{1}{\sqrt{40}}}{1 - \dfrac{4}{5\sqrt{40}}}
 = \frac{5}{6} + \frac{\sqrt{10}}{12}
\]
\hfill(30\%)

\solhead{3.6}{}
% source id: 2008_TH_S9
\ioaafig{2008_TH_S9-s1.pdf}
Power of radiation received by the Earth from the Sun is
\[
P_m = \sigma T_{\mathrm{eff}\,\odot}^{4}\, 4\pi R_\odot^2\,
      \frac{\pi\, r_\oplus^2}{4\pi\, x^2}\, t_1
\]
\hfill(40\%)

Power radiated by the Earth and reach the outer space
\[
P_{\mathrm{out}} = \sigma\, T_\oplus^4\, 4\pi\, r_\oplus^2\, t_2
\]
\hfill(30\%)

In equilibrium $P_{\mathrm{in}} = P_{\mathrm{out}}$ \hfill(5\%)
\begin{align*}
\sigma T_{\mathrm{eff}\,\odot}^4\,\frac{\pi r_\oplus^2}{4\pi x^2}\,
  4\pi R_\odot^2 t_1 &= \sigma T_\oplus^4\, 4\pi r_\oplus^2 t_2\\
T_\oplus^4 &= \frac{T_{\mathrm{eff}\,\odot}^4\, t_1 R_\odot^2}{4 t_2 x^2}\\
T_\oplus &= \sqrt[4]{\frac{T_{\mathrm{eff}\,\odot}^4\, t_1 R_\odot^2}
                          {4 t_2 x^2}}
          = T_{\mathrm{eff}\,\odot}
            \sqrt[4]{\frac{t_1}{4 t_2}\left(\frac{R_\odot}{x}\right)^2}
\end{align*}
\hfill(25\%)

\solhead{3.7}{}
% source id: 2009_TH_P13
\ioaafig{2009_TH_P13-s1.pdf}
We know the orbital period of Mars then the angle $\angle M_1 S M_2$ can be
determined simply
\[
\angle M_1 S M_2 = \frac{106}{687}\times 360 = 55.5^\circ
\]
\hfill(2 points)

By the same way we determine $\angle E_1 S E_2$:
\[
\angle E_1 S E_2 = \frac{106}{365}\times 360 = 104.5^\circ
\]
\hfill(2 points)

Then angle $\angle M_2 S E_2$ is:
\[
\angle M_2 S E_2 = 104.5 - 55.5 = 49.0^\circ
\]
\hfill(2 points)
\[
\cos(\angle M_2 S E_2) = \frac{SE_2}{SM_2}
\]
\hfill(2 points)
\[
r_{\mathrm{mars}} = \frac{SM_2}{SE_2}
 = \frac{1}{\cos(\angle M_2 S E_2)}
 = \frac{1}{\cos(49^\circ)} = 1.52\ \mathrm{AU}
\]
\hfill(2 points)

\solhead{3.8}{}
% source id: 2010_TH_L17
\begin{description}
\item[(a)] Drawing scaled diagram is impossible. Rough sketch is accepted.
\item[(b)] There are 10 days and nights per taris year. The obliquity is
  $3^\circ$, which means that the planet's rotation is in the same direction
  as its orbit. Thus, total number of rotations per year is $10 + 1 = 11$.

  \textbf{Note:} The obliquity is positive (similar to the Earth / Mars /
  Jupiter). This means, we have ADD one rotation. Subtracting one rotation by
  assuming opposite rotation (like the Venus) is incorrect.
\item[(c)] By Kepler's third law, $\dfrac{T^2}{R^3} = \text{Constant}$
  \begin{align*}
  \frac{T_{en}^2}{R_{en}^3} &= \frac{T_{ex}^2}{R_{ex}^3} \tag{1}\\
  R_{ex}^3 &= \frac{1.6^2 R_{en}^3}{0.2^2} \tag{2}\\
  R_{ex} &= \sqrt[3]{64}\ \text{endorlengths} \tag{3}\\
         &= 4\ \text{endorlengths} \tag{4}
  \end{align*}
\item[(d)] Using same logic as above
  \begin{align*}
  \frac{T_C^2}{R_C^3} &= \frac{T_T^2}{R_T^3} \tag{5}\\
  T_C^2 &= \frac{9^3 R_T^3 T_T^2}{R_T^3} \tag{6}\\
  T_C &= \sqrt{729}\ \text{tarisyears} \tag{7}\\
      &= 27\ \text{tarisyears} \tag{8}
  \end{align*}
\item[(e)] As Corulus is in Opposition, Sola - Taris - Corulus form straight
  line (in that order). Distance $= 9 - 1 = 8$ tarislengths.
\item[(f)] In the figure, S is Sola, A and B are start of the year positions
  of Taris and Corulus, T and C are their positions after `$n$' days. Angles
  are named from $a$ to $f$. The dashed line is parallel to line SB.
  Triangle(SCT) is used for sine rule as well as answer in the next part.
  Figure is not to the scale.
  \ioaafig{2010_TH_L17-s1.pdf}
  \begin{align*}
  a + b + c &= \pi \tag{9}\\
  b + d + e &= \pi \tag{10}\\
  d &= f + c \tag{11}\\
  f + c &= \frac{2\pi n}{10} \tag{12}\\
  f &= \frac{2\pi n}{270} \tag{13}\\
  \sin b &= 9\sin a \ \text{(By Sine Rule)} \tag{14}\\[6pt]
  e &= \pi - b - d \tag{15}\\
    &= \pi - b - c - f \tag{16}\\
    &= a - f \tag{17}\\[6pt]
  b &= \pi - (a + c) \tag{18}\\
    &= \pi - \left(a + \frac{2\pi n}{10} - \frac{2\pi n}{270}\right) \tag{19}\\
    &= \pi - \left(a + \frac{52\pi n}{270}\right) \tag{20}\\[6pt]
  9\sin a &= \sin\left(\pi - \left(a + \frac{52\pi n}{270}\right)\right)
    \tag{21}\\
    &= \sin\left(a + \frac{52\pi n}{270}\right) \tag{22}\\
    &= \left[\sin a\cos\left(\frac{52\pi n}{270}\right)
       + \cos a\sin\left(\frac{52\pi n}{270}\right)\right] \tag{23}\\
  9 &= \cos\left(\frac{52\pi n}{270}\right)
       + \cot a\sin\left(\frac{52\pi n}{270}\right) \tag{24}\\
  \cot a &= \frac{9 - \cos\left(\frac{52\pi n}{270}\right)}
                 {\sin\left(\frac{52\pi n}{270}\right)} \tag{25}\\
  a &= \tan^{-1}\!\left[\frac{\sin\left(\frac{52\pi n}{270}\right)}
       {9 - \cos\left(\frac{52\pi n}{270}\right)}\right] \tag{26}\\[6pt]
  \lambda &= \pi - e \tag{27}\\
          &= \pi + f - a \tag{28}\\
  \lambda &= \pi + \frac{2\pi n}{270}
     - \tan^{-1}\!\left[\frac{\sin\left(\frac{52\pi n}{270}\right)}
       {9 - \cos\left(\frac{52\pi n}{270}\right)}\right] \tag{29}
  \end{align*}
\item[(g)]
  \begin{align*}
  \text{Area} &= \frac12 \times l(ST) \times l(SC) \times \sin c\\
              &= \frac12 \times 1 \times 9 \times 0.568\\
              &= 2.56
  \end{align*}
  The area is about 3\,(tarislength)$^2$.
\end{description}

\solhead{3.9}{}
% source id: 2010_TH_S11
\[
T_M = \sqrt{\frac{R_M^3}{R_E^3}}\,T_E = 1.881\ \mathrm{years}
\]
\hfill(2 points)
\begin{align*}
\frac{1}{T_s} &= \frac{1}{T_E} - \frac{1}{T_M}\\
T_s &= \frac{T_E \times T_M}{(T_M - T_E)}
     = \frac{1.881}{0.881}\times 365.25 = 779.8\ \mathrm{days}
\end{align*}
\hfill(3 points)

That means there is an opposite of the Mars about every 780 days. If the next
great opposite will be in 2018, then
\begin{align*}
15\times 365 + 4 &= 5479\ \mathrm{days}\\
5479/779.8 &= 7.026
\end{align*}
It means that there will have been 7 opposites before Aug.\,28, 2018,
\hfill(3 points)

So the date for the great opposite should be
\[
5479 - 7\times 779.8 = 20.4\ \mathrm{days},
\]
i.e.\ 20.4 days before Aug.\,28, 2018, \hfill(2 points)

It is on Aug.\,7, 2018.

\solhead{3.10}{}
% source id: 2010_TH_S14
Since the observer is at the pole, the affect of the earth's rotation on the
transit could be neglected.

Then the Sun's angle at the earth extends as
\[
\theta_0 = \arcsin\!\left(\frac{2r_{\mathrm{sun}}}{1\,\mathrm{AU}}\right)
 \approx 32.0';
\]
the angular velocity of the Venus around the Sun, respected to the earth is
$\omega_1$,
\[
\omega_1 = \omega_{\mathrm{venus}} - \omega_{\mathrm{earth}}
 = \frac{2\pi}{T_{\mathrm{venus}}} - \frac{2\pi}{T_{\mathrm{earth}}}
 \approx 4.29\times 10^{-4}\ ('/\mathrm{s})
\]
\hfill(2 points)

For the observer on earth, Venus moved $\theta$ during the whole transit. Let
OE be perpendicular to AB,
\[
OA = 16', \quad \angle AOB = 90^\circ, \quad MN \parallel AB,
\]
So $OE = 11.3'$,
$OC = \dfrac{d'_{\mathrm{venus}}}{2} + \dfrac{r'_{\mathrm{sun}}}{2}$,
$d'_{\mathrm{venus}}$ is the angular size of Venus seen from Earth.
\ioaafig{2010_TH_S14-s1.pdf}
\ioaafig{2010_TH_S14-s2.pdf}
\begin{align*}
d'_{\mathrm{venus}} &= \frac{2\times 0.949\times 6378}
                           {(1-0.723)\times 1\,\mathrm{AU}} \approx 1',\\
OC &\approx 16.5', \quad CD \approx 24.0',\\
CE &= \sqrt{OC^2 - OE^2} \approx 12.0'\\
CD &= 2CE = 24.0'
\end{align*}
So, $\theta = \angle CFD = 24.0'$, \hfill(3 points)

As shown on the picture, $\theta' = \angle COD$ is the additional angle that
Venus covered during the transit,
\[
\frac{\operatorname{tg}\dfrac{\theta}{2}}{\operatorname{tg}\dfrac{\theta'}{2}}
 = \frac{0.723}{(1-0.723)}, \quad
\operatorname{tg}\frac{\theta}{2} = \operatorname{tg} 12', \quad
\theta' = 9.195';
\]
\hfill(3 points)
\[
t_{\mathrm{transit}} = \frac{\theta'}{\omega_1}
 = \frac{9.195'}{4.29\times 10^{-4}\ '/\mathrm{s}}\times\cos\varepsilon,
 \text{ that is } 5^{\mathrm{h}}56^{\mathrm{m}}36^{\mathrm{s}},
\]
% sic: the factor cos(epsilon) appears in the source without epsilon being
% defined anywhere in the solution.
So the transit will finish at about $14^{\mathrm{h}}57^{\mathrm{m}}$.
\hfill(2 points)

\solhead{3.11}{}
% source id: 2010_TH_S15
The semi-major axis of Moon's orbit is $a$; its eccentricity is $e$; $T$ is
the revolution period; apparent radius of the Moon is $r$; the distance
between Earth and Moon is $d$; the angular radius of the Sun is $R_\odot$.

When the Moon is at perigee, the total eclipse will be longest.
\[
\omega_1 = v_1/d_1, \qquad t_1 = 2(r_1 - R)/\omega_1
\]
Here, $\omega$ is the angular velocity of the moon, and $v$ is its linear
velocity; $t_2$ is the during time of total solar eclipse; $r_1$ is the
angular radius of the Moon when it's at perigee.
% sic: the source defines "t2" as the duration of the total eclipse although
% the formula above uses t1.

When the Moon is at apogee, the annular eclipse will be longest.
\[
\omega_2 = v_2/d_2, \qquad t_2 = 2(R - r_2)/\omega_2
\]
Since $v_2/v_1 = d_1/d_2 = (1-e)/(1+e)$, we get:
\begin{equation}
\frac{t_2}{t_1} = \frac{R-r_2}{r_1-R}\times\left(\frac{1+e}{1-e}\right)^2
\tag{1}
\end{equation}
\hfill(3 points)

Moon orbits the Earth in a ellipse. Its apparent size $r$ varies with time.
When $r > R$, if there occurred an center eclipse, it must be total solar
eclipse. Otherwise when $r \le R$, the center eclipse must be annular.
% sic: "in a ellipse", "an center eclipse" as in the original.

We need to know that, in a whole moon period, what's the time fraction of
$r > R$ and $r \le R$. $r \propto 1/d$. But it's not possible to get $d$ by
solving the Kepler's equation. Since $e$ is a small value, it would be
reasonable to assume that $d$ changes linearly with $t$. So, $r$ also changes
linearly with $t$. Let the moment when the Moon is at perigee be the starting
time ($t=0$), in half a period, we get:
\[
r = r_2 + kt = r_2 + \frac{2(r_1-r_2)}{T}\cdot t, \qquad 0 \le t < T/2
\]
Here, $k = 2(r_1-r_2)/T = \mathrm{constant}$.

When $r = R$, we get a critical $t$:
\begin{equation}
t_R = \frac{R-r_2}{k} = \frac{(R-r_2)}{2(r_1-r_2)}\cdot T
\tag{2}
\end{equation}
\hfill(2 points)

During a Moon period, if $t \in (t_R, T-t_R)$, then $r > R$, and the central
eclipses occurred are total solar eclipses. The time interval from $t_R$ to
$T-t_R$ is $\Delta t_T = T - 2t_R$. If $t \in [0, t_R]$ \&\ $t \in [T-t_R, T]$,
then $r \le R$, and the central eclipses occurred are annular eclipses. The
time interval is $\Delta t_A = 2t_R$.
\hfill(4 points)

The probability of occurring central eclipse at any $t$ is the same. Thus the
counts ratio of annular eclipse and total eclipse is:
\[
\frac{f_A}{f_T} = \frac{\Delta t_A}{\Delta t_T} = \frac{2t_R}{T-2t_R}
 = \frac{R-r_2}{r_1-R}
 = \left.\frac{t_2}{t_1}\right/\left(\frac{1+e}{1-e}\right)^2
 \approx \frac{4}{3}
\]
\hfill(1 point)

\solhead{3.12}{}
% source id: 2011_TH_S10
Formulation of a penetration conditions \hfill[2]

From the conditions specified in the exercise text the particle should be pass in the belt of
magnetic field around quarter of the circle. If this particle leaves the belt the particle path
should be a little longer. It means that the radius of this circle should be larger then the belt
thickness.
% sic: "should be pass", "larger then" as in the original.

Determination of the forces (centrifugal and magnetic) acting upon the particle, writing of the
appropriate formulae \hfill[4]

Comparison of the forces and calculations \hfill[3]

Comparing centrifugal and magnetic force one obtains
\[
e \cdot v \cdot B = M \cdot v^2 / r
\]
So $e \cdot r \cdot B = p$, where $p$ denotes the momentum. From the assumed approximation
$e \cdot r \cdot B = E / c$.

Thus definitively particle energy is described by $E = c \cdot e \cdot r \cdot B$.

Substituting numerical values one obtains $E \approx 3\cdot10^{-8}$ J it means
$2\cdot10^{2}$ GeV.

Answer and conclusion: \hfill[1]

\solhead{3.13}{}
% source id: 2012_TH_S3
\ioaafig{2012_TH_S3-s1.pdf}

From Kepler's Third Law, we have:
\[
\frac{a_{\mathrm{Mars}}^3}{T_{\mathrm{Mars}}^2} = \frac{a_\oplus^3}{T_\oplus^2}
\;\Rightarrow\; T_{\mathrm{Mars}} = 1.874\ year
% original uses the planetary symbol for Mars as the subscript
\]
\hfill(5 pt -- calculating the period of Mars)

Thus
\[
\frac{1}{T_{\mathrm{op}}} = \frac{1}{T_\oplus} - \frac{1}{T_{\mathrm{Mars}}}
\;\Rightarrow\; T_{\mathrm{op}} = 2.14\ year
\]
\hfill(5 pt -- synodic period)

\solhead{3.14}{}
% source id: 2012_TH_S4
The Moon's flux is proportional to its albedo, so $\frac{F_1}{a} = \frac{F_2}{1}$,
being $F_1$ its actual flux, $a$ the albedo and $F_2$ the hypothetical flux in
the case the Moon had an albedo 1. \hfill(5 pt)

Being $m_1$ the Moon apparent magnitude and $m_2$ the desired value, we have:
\begin{equation*}
m_2 - m_1 = -2.5\log\frac{F_2}{F_1}
\quad\Rightarrow\quad
m_2 = m_1 + 2.5\log a \tag*{(5 pt)}
\end{equation*}
Replacing $m_1 = -12.74$ and $a = 0.14$ we have: $m_2 = \mathbf{-14.9}$

\solhead{3.15}{}
% source id: 2013_TH_S1
\textit{Answer:}
\[
T_\oplus = T_\odot\,(1-\alpha)^{1/4}\sqrt{\frac{R_\odot}{2r_\oplus}}
         = 246~\mathrm{K}
\quad\text{or}\quad
T_\oplus = -27\,^\circ\mathrm{C}
\]
\hfill(10 Points)

(Note: $r_\oplus$ = average Earth--Sun distance; the albedo, $a = 0.39$, is
given in the additional material.)

\solhead{3.16}{Apparent magnitude of the Moon}
% source id: 2014_TH_P14
The apparent magnitude of a planet from the Solar System depends on the phase
angle $M = M(\Psi)$. The apparent magnitude of the body is given by the
relation:
\[
m = M + 2.5\cdot\log\frac{d_{C,S}^{2}\cdot d_{C,O}^{2}}{d_{0}^{4}\cdot p(\Psi)},
\]
% sic: subscripts C,S / C,O here vs B,S / B,O in the text below, as in the original
unde: $d_{B,S}$ -- the distance between the body and the Sun; $d_{B,O}$ --
distance between the body and observer; $d_{0} = 1$\,AU; $\Psi$ -- the phase
angle; $p(\Psi)$ -- the phase function:
% sic: "unde" (Romanian for "where") as in the original
\[
p(\Psi) = \frac{2}{3}\cdot\left[\left(1 - \frac{\Psi}{\pi}\right)\cos\Psi
  + \frac{1}{\pi}\sin\Psi\right],
\]
$\Psi$ as seen in the figure bellow is given by the cosine law.
% sic: "bellow" as in the original
\ioaafig{2014_TH_P14-s1.pdf}
\[
\cos\Psi = \frac{d_{BO}^{2} + d_{BS}^{2} - d_{OS}^{2}}{2 d_{BO}\cdot d_{BS}}.
\]
In particularaly for the Moon
% sic: "particularaly" as in the original
\ioaafig{2014_TH_P14-s2.pdf}
\begin{gather*}
\cos\Psi = \frac{d_{ME}^{2} + d_{MS}^{2} - d_{ES}^{2}}{2 d_{ME}\cdot d_{MS}};\\
p(\Psi) = \frac{2}{3}\cdot\left[\left(1 - \frac{\Psi}{\pi}\right)\cos\Psi
  + \frac{1}{\pi}\sin\Psi\right];\\
m_{M} = M_{M} + 2.5\cdot\log\frac{d_{M,S}^{2}\cdot d_{ME}^{2}}
  {d_{0}^{4}\cdot p(\Psi)}.
\end{gather*}

Particular cases:

1) Full moon
\begin{gather*}
\Psi = 0;\\
\cos\Psi = 1;\quad \sin\Psi = 0;\\
p(\Psi) = \frac{2}{3};\\
d_{MS} = 1~\mathrm{AU};\ d_{ME} = 385000~\mathrm{km} \approx 0.00256~\mathrm{AU}
  = 256\cdot10^{-5}~\mathrm{SU};\ d_{0} = 1~\mathrm{SU};
% sic: "SU" as in the original
\end{gather*}
\[
m_{M} = M_{M} - 12.5^{\mathrm{m}} = 0.25^{\mathrm{m}} - 12.5^{\mathrm{m}}
      = -12.25^{\mathrm{m}}.
\]

2) First Quarter
\begin{gather*}
\Psi = 90^\circ;\\
\cos\Psi = 0;\quad \sin\Psi = 1;\\
p(\Psi) = \frac{2}{3\pi} \approx 0.2;\\
\frac{d_{M,S}^{2}\cdot d_{ME}^{2}}{d_{0}^{4}\cdot p(\Psi)}
  = \frac{65536\cdot10^{-10}}{0.2} = 491520\cdot10^{-10};
% sic: 65536/0.2 = 327680, source prints 491520
\end{gather*}
\[
m_{M} = M_{M} - 10.75^{\mathrm{m}} = 0.25^{\mathrm{m}} - 10.75^{\mathrm{m}}
      = -10.5^{\mathrm{m}}.
\]

\solhead{3.17}{}
% source id: 2015_TH_S3
\begin{enumerate}
\item \textbf{Calculate the Earth revolution speed.} \hfill(20)

  Earth revolution speed is
  \[
  v_E = \frac{2\pi \bar r_E}{T_E} = \frac{2\pi a_E}{T_E} = 29805~\mathrm{m/s}
  \]
\item \textbf{Calculate Julietta revolution speed.} \hfill(25)

  Julietta revolution is (since the given data is semimajor axis, the orbit of
  Julietta is ellipse)
  \[
  v_J = \sqrt{GM\left(\frac{2}{r_J} - \frac{1}{a}\right)} = 16742.9~\mathrm{m/s}
  \]
\item \textbf{Calculate the approximate dimension of Julietta.}

  From the geometry
  \ioaafig{2015_TH_S3-s1.pdf}
  \[
  \cos\alpha = \frac{r_E^2 + r_J^2 - r_{J-E}^2}{2 r_E r_J} = 0.995
  \]
  % sic: (1 + 3.076^2 - 2.156^2)/(2 x 3.076) = 0.945; source prints 0.995,
  % while the numbers below are consistent with 0.945
  Thus, Julietta's relative projected speed to Earth is
  \[
  v_{J,\mathrm{rel}} = v_E - v_J\cos\alpha = 13984.16~\mathrm{m/s} \tag*{(40)}
  \]
  Hence, the approximate dimension of Julietta is
  \[
  l_J = v_{J,\mathrm{rel}}\,t_{\mathrm{occil}} = 86701.8~\mathrm{m} \tag*{(15)}
  \]
  % sic: "t_occil" as in the original; 13984.16 x 6.201 = 86715.6, source
  % prints 86701.8
\end{enumerate}

\solhead{3.18}{}
% source id: 2015_TH_S9
\begin{enumerate}
\item Solar luminosity is
  \begin{align*}
  L_\odot &= 4\pi R_\odot^2\sigma T_\odot^4\\
    &= 4\times3.14\times(6.96\times10^{10})^2\,(5.67\times10^{-5})(5780)^4~\mathrm{ergs/sec}\\
    &= 3.9\times10^{33}~\mathrm{ergs/sec} \tag*{(10)}
  \end{align*}
\item Flux of the Sun intercepted by Venus is
  \[
  L_V = \left(A_{V\mathrm{(projected)}}/A_{\mathrm{(sphere,0.72au)}}\right)\times L_\odot
  \]
  Also,
  \[
  L_V = 4\pi R_V^2\sigma T_V^4
  \]
  So,
  \[
  4\pi R_V^2\sigma T_V^4 = (\pi R_V^2/4\pi d^2)\times L_\odot
    = (\pi R_V^2/4\pi d^2)\times 4\pi R_\odot^2\sigma T_\odot^4 \tag*{(30)}
  \]
\item Therefore,
  \begin{align*}
  T_V^4 &= (R_\odot^2/4d^2)\times T_\odot^4
    = \left[(6.957\times10^5~\mathrm{km})^2/4\times(0.72\times149.6\times10^6)^2\right]\times(5780)^4\\
  &= (483998490000/46407499776000000)\times(5780)^4
    = (4.84\times10^{11}/4.64\times10^{16})\times(5780)^4\\
  &= 1.0429316216908190111241385226915\text{e-}5\times1114577187760656\\
  % sic: "e-5" calculator notation and unrounded digit strings as in the original
  &= 1.04\times10^{-5}\times1.11\times10^{15}\\
  &= 11624277939.308134327737156481455 = 1.1624\times10^{10}\\
  T_V &= 328.5~\mathrm{K} \tag*{(20)}
  \end{align*}
\item The flux density of Venus at closest approach to Earth at an observing
  frequency of 5\,GHz: ($1'' = 4.84\times10^{-6}$\,rad)
  \begin{align*}
  S_{\mathrm{5GHz}} &= B_{\mathrm{5GHz}}\,\Omega = (2kT_b\nu^2/c^2)\,\Omega
    \qquad(\rightarrow T_b = T_V)\\
  &= \left(2\times1.38\times10^{-16}\times328.5\times(5\times10^9)^2
    \times\pi\,(66\times4.84\times10^{-6})^2\right)/\,(3\times10^{10})^2\\
  % sic: the solid angle uses the full 66 arcsec angular diameter in place of
  % the 33 arcsec radius, as in the original
  &= 8.07\times10^{-22}~\mathrm{ergs~s^{-1}~cm^{-2}~Hz^{-1}}\\
  &= 80.7~\mathrm{Jy} \tag*{(40)}
  \end{align*}
\end{enumerate}

\solhead{3.19}{Gases on Titan}
% source id: 2016_TH_T2
As the gas is monatomic,
\begin{align*}
\frac{3}{2}k_{\mathrm{B}}T_{\mathrm{T}} &\approx \frac{1}{2}m_{\mathrm{g}}v_{\mathrm{rms}}^2\\
3k_{\mathrm{B}}T_{\mathrm{T}} &\approx \frac{M_{\mathrm{g}}}{N_{\mathrm{A}}}v_{\mathrm{rms}}^2\\
\therefore\ v_{\mathrm{rms}} &\approx \sqrt{\frac{3k_{\mathrm{B}}N_{\mathrm{A}}T_{\mathrm{T}}}{M_{\mathrm{g}}}}
\end{align*}
\hfill(4.0)

To remain in atmosphere,
\begin{align*}
v_{\mathrm{rms}} &< \frac{v_{\mathrm{esc}}}{6} = \frac{1}{6}\sqrt{\frac{2GM_{\mathrm{T}}}{R_{\mathrm{T}}}}\\
\sqrt{\frac{3k_{\mathrm{B}}N_{\mathrm{A}}T_{\mathrm{T}}}{M_{\mathrm{g}}}} &< \sqrt{\frac{GM_{\mathrm{T}}}{18R_{\mathrm{T}}}}\\
\therefore\ M_{\mathrm{g}} &> \frac{54k_{\mathrm{B}}N_{\mathrm{A}}T_{\mathrm{T}}R_{\mathrm{T}}}{GM_{\mathrm{T}}} \tag*{(4.0)}\\
&> \frac{54 \times 1.381\times10^{-23} \times 6.022\times10^{23} \times 93.7 \times 2.575\times10^{6}}
        {6.6741\times10^{-11} \times 1.23\times10^{23}}~\mathrm{g}\\
&> 13.2~\mathrm{g} \tag*{(1.5)}
\end{align*}

Thus, all gases with atomic weight more than $A_{\min} = 13.2$ will be
retained in the atmosphere of Titan. \hfill(0.5)

\medskip
\underline{Alternative solution}
\begin{align*}
\frac{3}{2}k_{\mathrm{B}}T_{\mathrm{T}} &\approx \frac{1}{2}m_{\mathrm{g}}v_{\mathrm{rms}}^2\\
\therefore\ v_{\mathrm{rms}} &\approx \sqrt{\frac{3k_{\mathrm{B}}T_{\mathrm{T}}}{m_{\mathrm{g}}}} \tag*{(4.0)}\\
\therefore\ m_{\mathrm{g}} &> \frac{54k_{\mathrm{B}}T_{\mathrm{T}}R_{\mathrm{T}}}{GM_{\mathrm{T}}} \tag*{(4.0)}\\
m_{\mathrm{g}} &> 2.19\times10^{-26}~\mathrm{kg} \tag*{(1.0)}\\
\therefore\ A_{\min} &= \frac{m_{\mathrm{g}}}{\text{atomic mass unit}}
  = \frac{2.19\times10^{-26}~\mathrm{kg}}{1.66\times10^{-27}~\mathrm{kg}}\\
A_{\min} &= 13.2 \tag*{(1.0)}
\end{align*}

\solhead{3.20}{The Star of Bethlehem}
% source id: 2017_TH_T8
\begin{description}
\item[a.] Assume that the average Earth's orbit is at the Sun. Therefore, the
  average great conjunction period equals to the synodic period between Jupiter
  and Saturn.
  \begin{align*}
  \frac{1}{P_{\mathrm{Synodic}}} &= \frac{1}{P_{\mathrm{Jupiter}}} - \frac{1}{P_{\mathrm{Saturn}}} \tag*{[3.0]}\\
  \frac{1}{P_{\mathrm{Synodic}}} &= \frac{1}{11.86} - \frac{1}{29.45}\\
  P_{\mathrm{Synodic}} &= 19.86~\text{years} \tag*{[1.0]}
  \end{align*}
  From average great conjunction period, in 19.86 years Jupiter and Saturn move
  \[
  \text{Heliocentric angle} = \frac{19.86}{11.86}\times360^\circ
    = 602.7^\circ = 242.7^\circ
  \]
  in an eastward direction through the zodiac, or
  $360^\circ - 242.7^\circ = 117.3^\circ$ in a westward direction through the
  zodiac. \hfill[2.0]
\item[b.] On 21st December 2020 (winter solstice), the Sun will be in the
  constellation of Sagittarius. Planets are always near the ecliptic. The
  current 12 zodiac constellations can be roughly divided into 12 equal parts
  of 30 degrees each.

  The conjunction will locate elongation $30.3^\circ$ East to the Sun.
  % sic: "will locate elongation" as in the original
  Therefore, the conjunction will be in the constellation of
  \textbf{Capricornus (Cap)}. \hfill[2.0]
\item[c.] The duration between 7BC and 2020 is 2026 years or 2027 years.

  Each conjunction occurs on the average every 19.86 years. Therefore, the
  number of conjunction between 7BC and 2020 is
  \[
  \text{Integer value of } 2026/19.86 = 102 \tag*{[2.0]}
  \]
  From a conjunction to its previous successive conjunction, the conjunction
  moves $243^\circ$ westwards. The position of the conjunction in 7BC is
  \[
  102\times242.7^\circ = 24758.33^\circ = 278.3^\circ \tag*{[3.0]}
  \]
  % sic: 102 x 242.7 = 24755.4, source prints 24758.33
  westward, or equivalently $81.7^\circ$ eastward from the position of the
  conjunction in 2020. Therefore, the conjunction was in the constellation of
  \textbf{Aries (Ari)} or \textbf{Pisces (Psc)}. \hfill[3.0]
\item[d.] At the second conjunction, Jupiter and Saturn had retrograde
  motions. For observers on the Earth, both of them were in opposition.
  Therefore, the sun was in the constellation of \textbf{Libra (Lib)} (if
  answer c) is Aries) or \textbf{Virgo (Vir)} (if answer c) is Pisces).
  \hfill[4.0]
\end{description}

\solhead{3.21}{Famous astronomical events}
% source id: 2019_TH_1
Correct solution: \hfill(10 p)

\begin{enumerate}
\item 1764: Discovery of the first planetary nebula
\item 1781: Discovery of Uranus
\item 1801: Discovery of Ceres
\item 1838: First successful measurement of a stellar parallax
\item 1877: Discovery of Phobos and Deimos
\item 1929: Discovery of the expansion of the Universe
\item 1943: Discovery of stellar populations
\item 1963: First identification of a quasar with an optical source
\item 1976: Viking probes arrived at planet Mars
\item 1986: Latest perihelion of comet 1/P Halley
\item 1990: Launch of the Hubble Space Telescope
\end{enumerate}

\solhead{3.22}{Flat Earth}
% source id: 2020_TH_2
Assume the surface area of \emph{one side} of the flat Earth is $A$. Let the angle between
the Sun and the flat Earth's center axis be $\theta$, where $\theta$ is initially $0^\circ$. As
the Sun's rays are parallel, the power delivered to the Earth by the Sun will be
$S_\odot A\cos\theta$ at any given point in time.

At equilibrium, this is the energy radiated away via blackbody radiation, so the equilibrium
temperature $T$ satisfies
\[
S_\odot A\cos\theta = \sigma(2A)T^4
\]
\hfill(2.0)

\noindent, where the factor of 2 comes from the fact that the flat Earth would radiate energy
from both sides.

This yields
\[
T(\theta) = \sqrt[4]{\frac{S_\odot\cos\theta}{2\sigma}}
\]
and we wish to find the value of $\theta_1$ such that $T(\theta_1) = T(0) - \Delta T$.
Thus,
\begin{align*}
\sqrt[4]{\frac{S_\odot\cos\theta_1}{2\sigma}} &= \sqrt[4]{\frac{S_\odot}{2\sigma}} - \Delta T
  & &\hspace{-2em}(2.0)\\
\cos\theta_1 &= \left(1 - \Delta T\sqrt[4]{\frac{2\sigma}{S_\odot}}\right)^4
  & &\hspace{-2em}(1.0)\\
 &= \left(1 - \sqrt[4]{\frac{2\times5.67\times10^{-8}}{1366}}\right)^4\\
\cos\theta_1 &= 0.9880
  & &\hspace{-2em}(1.0)
\end{align*}

Now, we find the time it takes for the axis to make such an angle with the Sun. On the
celestial sphere, let O be the center of precession, Z be the current direction of the axis,
and X be direction of the axis when it makes an angle of $\theta_1$ with the sun, so
$\angle ZCX = \theta_1$. If $\epsilon$ is the radius of precession, then
$\angle OCZ = \angle OCX = \epsilon$. \hfill(1.0)

By the spherical Law of Cosines on angle O of spherical triangle OXZ, we have
\begin{align*}
\cos\theta_1 &= \cos\epsilon\cos\epsilon + \sin\epsilon\sin\epsilon\cos(\sphericalangle O)\\
 &= \cos^2\epsilon + \sin^2\epsilon\cos(\sphericalangle O)\\
\sphericalangle O &= \cos^{-1}\left(\frac{\cos\theta_1 - \cos^2\epsilon}{\sin^2\epsilon}\right)
  & &\hspace{-2em}(1.0)\\
\therefore \Delta t = \frac{\sphericalangle O}{2\pi}\times P
 &= \frac{P}{2\pi}\times\cos^{-1}\left(\frac{\cos\theta_1 - \cos^2\epsilon}{\sin^2\epsilon}\right)
  & &\hspace{-2em}(1.0)\\
 &= \frac{25800}{2\pi}\times\cos^{-1}\left(\frac{0.9880 - \cos^2 23.5^\circ}{\sin^2 23.5^\circ}\right)\\
 &\approx 1606\,\mathrm{yr}
\end{align*}

Thus, the average temperature of the earth will go down by $1\,^\circ\mathrm{C}$ in just over
\textbf{1600 years}. \hfill(1.0)

\solhead{3.23}{Jupiter's Great Red Spot}
% source id: 2020_TH_8
\begin{description}
\item[(a)] From the figure, the centre of GRS lies at $\phi = -20.0^\circ$. \hfill(1.0)
\begin{align*}
\epsilon &= \frac{R_e}{R_p} = \frac{7.15}{6.69}\\
 &= 1.069 & &\hspace{-2em}(1.0)\\
\overline{r} = r(-20^\circ) &= R_e\left(1 + \epsilon^{-2}\tan^2\phi\right)^{-1/2}\\
 &= 7.15\times10^7 \times \left(1 + (1.0688)^{-2}\tan^2(-20^\circ)\right)^{-1/2}\\
\overline{r} &= 6.77\times10^7\,\mathrm{m} & &\hspace{-2em}(1.0)\\
\overline{R} = R(-20^\circ) &= R_e\epsilon^{-2}\left(\frac{r(\phi)}{R_e\cos\phi}\right)^3\\
 &= 7.15\times10^7 \times (1.0688)^{-2}
   \left(\frac{r(-20^\circ)}{7.15\times10^7\times\cos(-20^\circ)}\right)^3\\
\overline{R} &= 6.40\times10^7\,\mathrm{m} & &\hspace{-2em}(1.0)
\end{align*}

\item[(b)] We use the figure to estimate size major and minor axis of the GRS and hence
estimate the eccentricity.
\begin{align*}
a &= \overline{r}\theta_{long}\\
 &= 6.77\times10^7 \times \frac{8.2^\circ\times\pi}{180^\circ}\\
a &= 9.69\times10^6\,\mathrm{m} & &\hspace{-2em}(2.0)\\
b &= \overline{R}\theta_{lat}\\
 &= 6.40\times10^7 \times \frac{4.3^\circ\times\pi}{180^\circ}\\
b &= 4.80\times10^6\,\mathrm{m} & &\hspace{-2em}(2.0)\\
e_{GRS} &= \sqrt{1 - \left(\frac{b}{a}\right)^2} = \sqrt{1 - \left(\frac{4.8}{9.7}\right)^2}\\
e_{GRS} &\approx 0.87 & &\hspace{-2em}(1.0)
\end{align*}

\item[(c)] The maximum value of wind velocity ($V_w$) is about $120\,\mathrm{m/s}$. \hfill(1.0)
\begin{align*}
A_{GRS} &= \pi a b & &\hspace{-2em}(1.0)\\
L_{GRS} &\approx \frac{2\pi a + 2\pi b}{2} = \pi a(1 + \sqrt{1 - e^2}) & &\hspace{-2em}(2.0)\\
\therefore \xi &= \frac{V_w L_{GRS}}{A_{GRS}} = \frac{V_w\pi a(1 + \sqrt{1 - e^2})}{\pi a b}\\
 &= \frac{120\times(1 + \sqrt{1 - 0.87^2})}{4.8\times10^6}\\
\xi &\approx 3.7\times10^{-5}\,\mathrm{s}^{-1} & &\hspace{-2em}(2.0)
\end{align*}

\item[(d)] The absolute vorticity will be given by
\begin{align*}
\xi_a &= \xi + 2\Omega\sin\phi\\
 &= \xi + \frac{2\times2\times\pi\sin-20^\circ}{P}
   = 3.7\times10^{-5} + \frac{4\pi\sin-20^\circ}{3.57\times10^4}\\
 &= 3.7\times10^{-5} - 1.2\times10^{-4}\\
\xi_a &= -8.3\times10^{-5} & &\hspace{-2em}(2.0)
\end{align*}

\item[(e)] The GRS is a cyclonic system. \hfill(1.0)

\item[(f)] For this, we have to write $\xi_a$ as a function of $\phi$ and then find root of the
function closest to the current latitude.
\begin{align*}
\xi_a = 0 &= \xi + 2\Omega\sin\phi & &\hspace{-2em}(1.0)\\
0 &= \frac{V_w(1 + \sqrt{1 - e^2})}{b} + 2\Omega\sin\phi_1\\
0 &= \frac{V_w(1 + \sqrt{1 - e^2})}{\overline{R}_1\theta_{lat}} + 2\Omega\sin\phi_1\\
0 &= \frac{V_w\epsilon^2 R_e^2(1 + \sqrt{1 - e^2})\cos^3\phi}{\theta_{lat}r^3(\phi_1)}
   + 2\Omega\sin\phi_1\\
0 &= \frac{V_w\epsilon^2 R_e^2(1 + \sqrt{1 - e^2})\cos^3\phi_1}
   {\theta_{lat}(R_e\left(1 + \epsilon^{-2}\tan^2\phi_1\right)^{-1/2})^3}
   + 2\Omega\sin\phi_1 & &\hspace{-2em}(3.0)\\
0 &= \left(\frac{V_w\epsilon^2(1 + \sqrt{1 - e^2})}{\theta_{lat}R_e}\right)
   \left(\cos\phi_1\sqrt{1 + \epsilon^{-2}\tan^2\phi_1}\right)^3 + 2\Omega\sin\phi_1\\
-3.520\times10^{-4}\sin\phi_1 &= 3.81\times10^{-5}\times
   \left(1 - (1 - \epsilon^{-2})\sin^2\phi_1\right)^{3/2}\\
35.2^2\sin^2\phi_1 &= 3.81^2\times\left(1 - (1 - \epsilon^{-2})\sin^2\phi_1\right)^3
\end{align*}
Let us define $\sin^2\phi_1 = x$
\[
\therefore x = 0.01172\,(1 - 0.1245x)^3
\]
\hfill(4.0)

Iterating over $x$, with starting value as $x_0 = \sin^2(-20^\circ) = 0.11698$
\begin{center}
\begin{tabular}{cc}
$\mathbf{x_0}$ & $\mathbf{x_1}$\\
\hline
0.11698 & 0.01121\\
0.01121 & 0.01167\\
0.01167 & 0.01167\\
\end{tabular}
\end{center}
\hfill(2.0)

Thus,
\begin{align*}
\phi_1 &= -\sin^{-1}(\sqrt{0.01167})\\
\phi_1 &= -6.2^\circ \approx -6^\circ & &\hspace{-2em}(2.0)
\end{align*}

An alternative solution is to try to find this using the plotting method, but one needs to
make sure the result is obtained to the equivalent precision.
\end{description}

\solhead{3.24}{Temperature of the Earth}
% source id: 2021_TH_Q2
\begin{description}
\item[2.1] The solar radiation intensity ($I_T$) received on Earth is:
  \[
  I_T = \frac{P_s}{4\pi r_{s-t}^2} = \sigma T_s^4\cdot\left(\frac{R_s}{r_{s-t}}\right)^{2} \tag*{[1 point]}
  \]
  with $P_s = 4\pi\sigma R_s^2 T_s^4$.

  Also, Earth would absorb energy at this rate:
  \[
  P_{abs} = I_T\,\pi R_t^2 = \pi\sigma T_s^4\cdot\left(\frac{R_s\cdot R_t}{r_{s-t}}\right)^{2}
  \]
  With $R_t = 6.4\times10^{6}$\,m as the radius of the planet ``disk''.
  \hfill[1 point]

  Then, by thermal equilibrium, the absorbed radiation would be radiated over
  the planet's surface:
  \[
  P_{abs} = P_{rad}
  \]
  With
  \begin{gather*}
  P_{rad} = 4\pi\sigma R_t^2 T_t^4\\
  \pi\sigma T_s^4\cdot\left(\frac{R_s\cdot R_t}{r_{s-t}}\right)^{2}
    = 4\pi\sigma R_t^2 T_t^4 \tag*{[1 point]}\\
  T_t = T_s\cdot\left(\frac{R_s}{2\,r_{s-t}}\right)^{\frac{1}{2}}
    = 278.58~K = 5.43~^\circ C \tag*{[1 point]}
  \end{gather*}
  This would be very cold but still viable to harbor life.
\item[2.2] The Earth's absorbed radiation, considering the albedo, is:
  \begin{gather*}
  P'_{abs} = 0.7\cdot\pi\sigma T_s^4\cdot\left(\frac{R_s\cdot R_t}{r_{s-t}}\right)^{2} \tag*{[1 point]}\\
  T'_T = T_s\,(0.7)^{\frac{1}{4}}\cdot\left(\frac{R_s}{2\,r_{s-t}}\right)^{\frac{1}{2}}
    = 254.81~K = -18.34~^\circ C \tag*{[1 point]}
  \end{gather*}
\item[2.3] If Earth reabsorbs 58\% of the 70\% re emitted energy, then:
  \[
  T''_T = T_s\,[0.7 + (0.58\cdot0.7)]^{\frac{1}{4}}
    \cdot\left(\frac{R_s}{2\,r_{s-t}}\right)^{\frac{1}{2}}
    = 285.68~K = 12.53~^\circ C \tag*{[4 points]}
  \]
\end{description}

\solhead{3.25}{Solar Retrograde Motion on Mercury}
% source id: 2022_TH_9
From Kepler's 3\textsuperscript{rd} Law, $T^2 = a^3$ in units of Earth years and au respectively.
\[
\therefore T_{\mathrm{Mercury}} = \sqrt{(0.387)^3} = 0.240\,75\,\mathrm{yr} = 87.934\,\mathrm{days}
\]
\hfill(2.0)

We know that $\omega_{\mathrm{orb}} = v/r$ so will vary throughout the orbit (since $r$ varies within
an ellipse) whilst $\omega_{\mathrm{rot}}$ is constant.
\[
\omega_{\mathrm{rot}} = \frac{2\pi}{T_{\mathrm{rot}}}
= \frac{2\pi}{\frac{2}{3} \times 87.934 \times 86400}
= 1.24 \times 10^{-6}\,\mathrm{rad/s}
\]
\hfill(2.0)

For retrograde motion, we need $\omega_{\mathrm{orb}} \ge \omega_{\mathrm{rot}}$.\\
Using the vis-viva equation for the critical value of $r$ for when
$\omega_{\mathrm{rot}} = \omega_{\mathrm{orb}}$,
\begin{align*}
\omega_{\mathrm{orb}}^2 &= \frac{v^2}{r^2}
  = \frac{GM_\odot}{r^2}\left(\frac{2}{r} - \frac{1}{a}\right) = \omega_{\mathrm{rot}}^2\\
\frac{2GM_\odot}{r^3} - \frac{GM_\odot}{ar^2} &= \omega_{\mathrm{rot}}^2\\
\therefore \frac{\omega_{\mathrm{rot}}^2}{GM_\odot}r^3 + \frac{1}{a}r - 2 &= 0\\
(1.160\times10^{-32})r^3 + (1.725\times10^{-11})r - 2 &= 0 \quad
  \text{(if $r$ is in meters)}\\
38.99r^3 + 2.584r - 2 &= 0 \quad \text{(if $r$ is in au)}
\end{align*}
\hfill(4.0)

Solving the cubic equation (by any valid method) gives only one non-imaginary root.
One possible way would be to use iterations for this polynomial $f(r)$, using perihelion
distance as the starting guess.
\begin{center}
\begin{tabular}{c|c}
$r$ & $f(r)$\\
\hline
0.3073 & $-0.07482$\\
0.3100 & $-0.03747$\\
0.3120 & $-0.00967$\\
0.3140 & 0.01841\\
0.3130 & 0.00433\\
0.3127 & 0.00012\\
\end{tabular}
\end{center}
\[
\therefore r = 0.3127\,\mathrm{au} = 4.684\times10^{10}\,\mathrm{m}
\]
\hfill(4.0)

From knowledge of ellipses, if $E$ is the eccentric anomaly
\begin{align*}
r &= a(1 - e\cos E)\\
\therefore E &= \cos^{-1}\left[\frac{1}{e}\left(1 - \frac{r}{a}\right)\right]
  = \cos^{-1}\left[\frac{1}{0.206}\left(1 - \frac{0.3127}{0.387}\right)\right]\\
E &= 0.3706\,\mathrm{rad} = 21.23^\circ = 21^\circ 14'
\end{align*}
\hfill(3.0)

Using Kepler's Equation we can find the mean anomaly, $M$
\begin{align*}
M &= E - e\sin E = 0.3706 - 0.206 \times \sin 0.3706\\
  &= 0.2960\,\mathrm{rad} = 16.96^\circ = 16^\circ 58'
\end{align*}
\hfill(2.0)

This is relative to the perihelion, so symmetry demands that the total time the sun is in
retrograde corresponds to
\begin{align*}
\Delta M &= 2M\\
\therefore T_{\odot,retro} &= T_{\mathrm{Mercury}}\frac{\Delta M}{2\pi}
  = 87.934 \times \frac{2 \times 0.296}{2\pi}\\
T_{\odot,retro} &= \boxed{8.28\,\mathrm{days}}
\end{align*}
\hfill(3.0)

[Accept alternative methods making use of the true anomaly e.g.\ with the relation
$\cos E = \frac{e + \cos\nu}{1 + e\cos\nu}$.]

\solhead{3.26}{Europa}
% source id: 2023_TH_Q4
\begin{description}
\item[(a)] From the formula $Q = \lambda S \Delta T t / d$, calculate the power
  $P$ of the heat flowing through a unit of the surface of ice:
  \[
    P = Q/t = \lambda S \Delta T / d.
  \]
  \hfill(1 point)

  Approximating the ice crust of Europa as a `wall', i.e.\ that the upper and
  lower surfaces are of equal area $S$, we obtain the power per unit surface
  area:
  \[
    P/S = \lambda \Delta T / d.
  \]
  \hfill(1 point)

  Substituting the data we obtain
  \[
    P/S = 86.5 \times 10^{-3}\ \mathrm{W\,m^{-2}}
        \approx 87\ \mathrm{mW\,m^{-2}},
  \]
  \hfill(1 point)

  similar to the value given for the Earth.

  The total power emitted inside Europa is therefore equal to
  \[
    4\pi R_{\mathrm{Eu}}^2 \, (P/S) = 2.65 \times 10^{12}\ \mathrm{W}.
  \]
  \hfill(2 points)

\item[(b)] The total power emitted inside the Earth is equal to
  \[
    4\pi R_{\oplus}^2 \times 70 \times 10^{-3}\ \mathrm{W\,m^{-2}}
      = 36 \times 10^{12}\ \mathrm{W},
  \]
  \hfill(1 point)

  which is about $13.5\times$ larger than the value for Europa. However, the
  mass ratio is 124. For the heat in Europa's interior to have a purely
  radioactive origin, the matter on Europa would have to contain an order of
  magnitude more radioactive elements per unit mass, which excludes this
  explanation.
  \hfill(3 points)

  Thus the answer must be tidal forces.
  \hfill(1 point)
\end{description}

\solhead{3.27}{Libration}
% source id: 2023_TH_Q7
\begin{description}
\item[(a)] $\phi_B = \alpha = 6^\circ 41'$
  \hfill(\textit{conceptually difficult:} 5 points)

\ioaafig{2023_TH_Q7-s1.pdf}

\item[(b)] The maximum angle of libration $\phi_L$ occurs when the Moon is at
  point $M$ and for $e = 0.064$:

\ioaafig{2023_TH_Q7-s2.pdf}

  $\phi_L$ is the Earth--Moon--F2 angle; $|EF2| = 2c$.\\
  Therefore
  $\sin(\phi_L/2) = c/a = e \implies
   \phi_L = 2\arcsin(0.064) = 7.34^\circ = 7^\circ 20'$
  \hfill(5 points)

\item[(c)] In the less accurate approximation,
  \[
    (2\phi_B + 2\phi_L)/180^\circ = x/50\% \implies x = 8\%
  \]
  and the total area visible is $S \approx 50\% + x = 58\%$.\\
  This is less accurate as it counts the overlapping areas made visible by both
  librations twice.

  A slightly more accurate approximation is given by taking the average of
  $\phi_B$ and $\phi_L$, $7^\circ 01'$. The belt made visible by libration will
  thus be $2\pi R_{\mathrm{Moon}} \cdot 7^\circ/360^\circ \approx 213$ km wide.
  The surface area of the belt is then approximately
  $213\ \mathrm{km} \times 2\pi R_{\mathrm{Moon}}
   = 2.32 \times 10^{6}\ \mathrm{km}^2$, which accounts for about 6\% of the
  lunar surface. If the students does this algebraically before substituting,
  they also get: % sic: "the students does"
  \[
    \frac{2\pi R_{\mathrm{Moon}} \cdot 2\pi R_{\mathrm{Moon}}}
         {4\pi \, 2 R_{\mathrm{Moon}}^2}
    \cdot \frac{7^\circ}{360^\circ}
    = \pi \, \frac{7^\circ}{360^\circ} \approx 6\%.
  \]
  In either case the total area visible is $S \approx 50\% + 6\% = 56\%$.
  \hfill(less accurate solution 4 points, better solution 5 points)

\item[(d)] It is necessary to find a common multiple of the period of libration
  in latitude $P_B = 27.2122$\,d (the draconic month) and the period of
  libration in longitude $P_L = 27.5545$\,d (the anomalistic month).

  The simplest way is analogous to synodic periods:
  \[
    \frac{1}{T_{\mathrm{syn}}} = \frac{1}{T_1} - \frac{1}{T_2}
  \]
  thus
  \[
    \frac{1}{T} = \frac{1}{P_B} - \frac{1}{P_L}
    \implies T \approx 2190.5\,\mathrm{d} = 74\,\mathrm{lunations}.
  \]
  Or using continued fractions:
  \[
    \frac{P_L}{P_B} = \frac{275545}{272122}
    \approx 1 + \cfrac{1}{79 + \cfrac{1}{2 + \cfrac{1}{131 + \cfrac{1}{6 + \cdots}}}}\,.
  \]
  With the first term we approximate $P_L/P_B \approx 80/79$; since
  $80 \times 27.2122 \approx 2176.9$ d and $79 \times 27.5545 \approx 2176.8$ d
  this is already close enough, and $T$ is again equal to about 74 lunations.
  \hfill(3 points)

  Since we only need half the cycle to see everything once (as the Moon swings
  away and back again during the whole cycle) therefore
  $T/2 = 37\ \mathrm{months} \approx 3$ years is enough to see all of the
  potentially visible parts of the Moon's surface.
  \hfill(2 points)
\end{description}

\solhead{3.28}{Greatest Eclipse}
% source id: 2024_TH_T10
\begin{description}
\item[(a)] The declination is simply the complement of the polar angle:
  \begin{align*}
    \delta &= 90^\circ - \theta\\
    \therefore \delta_\odot &= 21^\circ 30' 15.9''\\
    \delta_{\mathrm{Moon}} &= 21^\circ 12' 18.4''
  \end{align*}
  \hfill(1.0 + 1.0 + 1.0)

\item[(b)] The local sidereal time at Greenwich corresponds by definition to
  the meridian of right ascension that is right above Greenwich. Since
  Greenwich corresponds to an azimuthal angle of $0^h$, the right ascension at
  any given point corresponds simply to the sum of the local sidereal time at
  Greenwich and the azimuthal angle:
  \begin{align*}
    \alpha &= \varphi + LST_{\mathrm{Greenwich}}\\
    \therefore \alpha_\odot &= 4^h 21^m 7.3^s\\
    \alpha_{\mathrm{Moon}} &= 4^h 21^m 12.6^s
  \end{align*}
  \hfill(1.0 + 1.0 + 1.0)

\item[(c)] By subtracting the Cartesian coordinates of the Sun from the
  Cartesian coordinates of the Moon, it is possible to obtain a vector for the
  direction of the axis of the Moon's shadow cone:
  \hfill(1.0)
  \begin{align*}
    \vec{u} &= \langle -1.339 \times 10^{11},\ 4.317 \times 10^{10},\
                       -5.544 \times 10^{10} \rangle\ \mathrm{m}\\
    |\vec{u}| &= 1.512 \times 10^{11}\ \mathrm{m}\\
    \therefore \hat{u} &= \langle -0.8855,\ 0.2855,\ -0.3666 \rangle
  \end{align*}
  \hfill(1.0 + 1.0 + 1.0)

\item[(d)] If one draws a vector in the same direction as $\hat{u}$ up to the
  earth's surface, then its magnitude can be found by the formula, where
  $\vec{M}$ is the position vector of the Moon, $k$ is a constant, and
  $R_\oplus$ is the radius of the Earth:
  \[
    |\vec{M} + k\hat{u}|^2 = R_\oplus^2
  \]
  \hfill(5.0)

  Solving for $k$:
  \begin{align*}
    0 &= k^2 (\hat{u} \cdot \hat{u}) + 2k(\hat{u} \cdot \vec{M})
         + \vec{M} \cdot \vec{M} - R_\oplus^2\\
      &= k^2 + 2k(\hat{u} \cdot \vec{M}) + |\vec{M}|^2 - R_\oplus^2
  \end{align*}
  Now, $|\vec{M}|^2 - R_\oplus^2 = 1.288 \times 10^{17}\ \mathrm{m}$ % sic: units given as m in the source
  \begin{align*}
    2(\hat{u} \cdot \vec{M}) &= -7.178 \times 10^{8}\ \mathrm{m}\\
    \therefore 0 &= k^2 - k(7.178 \times 10^{8}) + (1.288 \times 10^{17})\\
    k &= \frac{7.178 \pm \sqrt{(-7.178)^2 - 4 \times 12.88}}{2} \times 10^{8}
  \end{align*}
  The vector will intersect the earth's surface at two points. Hence two
  solutions. However, only the first intersection with the sphere is relevant
  in this case, so only the smallest solution to the equation should be
  considered:
  \begin{align*}
    k &= \frac{7.178 - \sqrt{(-7.178)^2 - 4 \times 12.88}}{2} \times 10^{8}\\
      &= 3.528 \times 10^{8}
  \end{align*}
  \hfill(4.0)

  The vector of location of the greatest eclipse corresponds to the position
  vector of the Moon plus $k$ times the direction vector $\hat{u}$:
  \begin{align*}
    \vec{p} &= \vec{M} + k\hat{u}\\
      &= \langle 6.091 \times 10^{6},\ -1.828 \times 10^{6},\
                 4.854 \times 10^{5} \rangle
  \end{align*}
  \hfill(2.0)

  It is possible to use the following expressions to convert this vector to
  the spherical system:
  \begin{align*}
    r &= \sqrt{x^2 + y^2 + z^2}\\
      &= 6.378 \times 10^{6}\ \mathrm{m}\\
    \theta &= \arctan\left(\frac{\sqrt{x^2 + y^2}}{z}\right)
      \quad \text{(valid for } z > 0\text{)}\\
      &= 1.495\ \mathrm{rad} = 85^\circ 38'\\
    \varphi &= \arctan\left(\frac{y}{x}\right)
      \quad \text{(valid for } x > 0\text{)}\\
      &= -0.2915\ \mathrm{rad} = -16^\circ 42'
  \end{align*}
  \hfill(2.0)

  The latitude corresponds to the complement of the angle $\theta$, and the
  longitude corresponds to the angle $\varphi$:
  \[
    \boxed{4^\circ 22'\,\mathrm{N},\ 16^\circ 42'\,\mathrm{W}}
  \]
  \hfill(2.0)

\item[(e)] The approximate geometry described in the problem statement is
  shown below:

\ioaafig{2024_TH_T10-s1.pdf}
  \hfill(3.0)

  The angle $\beta$ can be calculated as follows:
  \begin{align*}
    \beta &= \arccos\left(
        \frac{R_\odot - R_{\mathrm{Moon}}}{d_\odot - d_{\mathrm{Moon}}}\right)\\
      &= \arccos\left(
        \frac{6.955 \times 10^{8} - 1.737 \times 10^{6}}
             {1.516 \times 10^{11} - 3.589 \times 10^{8}}\right)\\
      &= 89.74^\circ
  \end{align*}
  \hfill(3.0)

  Using this angle, it is possible to determine the radius of the umbra:
  \[
    r_{\mathrm{umbra}} = \frac{R_{\mathrm{Moon}}
      - (d_{\mathrm{Moon}} - R_\oplus)\cdot\cos(\beta)}{\sin(\beta)}
  \]
  \hfill(3.0)
  \[
    \boxed{r_{\mathrm{umbra}} = 1.196 \times 10^{5}\ \mathrm{m}}
  \]
  \hfill(1.0)

\item[(f)] The rotational velocity of the Earth at the latitude of the center
  of the umbra is the following:
  \begin{align*}
    v_{\mathrm{rot}} &= \frac{2\pi R_\oplus}{23^h 56^m 04^s}
        \cos(\phi_{\mathrm{umbra}})\\
      &= \frac{2\pi \times 6.378 \times 10^{6}}{23^h 56^m 04^s}
        \cos(4^\circ 22')
  \end{align*}
  \hfill(2.0)
  \[
    \boxed{v_{\mathrm{rot}} = 463.7\ \mathrm{m/s}}
  \]
  \hfill(1.0)

\item[(g)] Using the vis-viva equation:
  \begin{align*}
    v_{\mathrm{Moon}} &= \sqrt{GM_\oplus\left(\frac{2}{d_{\mathrm{Moon}}}
        - \frac{1}{a_{\mathrm{Moon}}}\right)}\\
      &= \sqrt{6.674 \times 10^{-11} \times 5.972 \times 10^{24}
        \left(\frac{2}{3.589 \times 10^{8}}
            - \frac{1}{3.844 \times 10^{8}}\right)}
  \end{align*}
  \hfill(2.0)
  \[
    \boxed{v_{\mathrm{Moon}} = 1088\ \mathrm{m/s}}
  \]
  \hfill(2.0)

\item[(h)] It is possible to use the following diagram to determine the value
  of the angle $\kappa$ between the $x$-axis and the velocity vector of the
  Moon:

\ioaafig{2024_TH_T10-s2.pdf}
  \hfill(6.0)

  Note that the great circle that intersects the Moon's meridian of right
  ascension at a right angle has a tangent line that coincides with the
  $x$-axis, so the angle $\kappa$ on the figure is also the angle between the
  velocity vector of the Moon and the $x$-axis. Also note that the Moon is
  moving from west to east in the celestial sphere (to the left in the
  figure), so its velocity vector is on the first quadrant of the $xy$-plane
  on this system.

  Using the four parts formula:
  \begin{align*}
    \cos(\delta_{\mathrm{Moon}})\cos(90^\circ)
      &= \sin(\delta_{\mathrm{Moon}})
         \cot(\alpha_{\mathrm{Moon}} + 24^h - \alpha_{\mathrm{intersection}})\\
      &\quad - \sin(90^\circ)\cot(90^\circ - \kappa)
  \end{align*}
  \hfill(3.0)
  \begin{align*}
    \therefore 0 &= \sin(\delta_{\mathrm{Moon}})
        \cot(\alpha_{\mathrm{Moon}} - \alpha_{\mathrm{intersection}})
        - \tan(\kappa)\\
    \tan(\kappa) &= \frac{\sin(\delta_{\mathrm{Moon}})}
        {\tan(\alpha_{\mathrm{Moon}} - \alpha_{\mathrm{intersection}})}\\
    \kappa &= \arctan\left[\frac{\sin(\delta_{\mathrm{Moon}})}
        {\tan(\alpha_{\mathrm{Moon}} - \alpha_{\mathrm{intersection}})}\right]\\
      &= \arctan\left[\frac{\sin(21^\circ 12' 18.4'')}
        {\tan(4^h 21^m 12.6^s - 23^h 07^m 59.2^s)}\right]\\
    \kappa &= 4^\circ 17'
  \end{align*}
  \hfill(2.0)

  Now it is possible to break down the velocity of the Moon into the $x$ and
  $y$ components:
  \[
    \mathbf{v}_{\mathrm{II}} =
    \begin{bmatrix}
      v_{\mathrm{Moon}}\cos(\kappa)\\
      v_{\mathrm{Moon}}\sin(\kappa)\\
      0
    \end{bmatrix}
    =
    \begin{bmatrix}
      1085\ \mathrm{m/s}\\
      81.25\ \mathrm{m/s}\\
      0
    \end{bmatrix}
  \]
  \hfill(3.0)

  The radial component of the velocity was neglected, so the velocity on the
  $z$-axis is equal to zero.

\item[(i)] It is possible to obtain system III by rotating system II from
  north to south by an angle of
  $\theta_{\mathrm{umbra}} - \theta_{\mathrm{Moon}}$ around the $x$-axis. The
  angles $\theta_{\mathrm{umbra}}$ and $\theta_{\mathrm{Moon}}$ correspond to
  the polar angles of the umbra and the Moon on system I.

  The following two figures illustrate this rotation.

\ioaafig{2024_TH_T10-s3.pdf}
  \hfill(3.0)

  Since the rotation is around the $x$-axis, the $x$ component remains
  unchanged. The component that was originally in the $y$-axis is broken down
  into two components on the $y$-axis and the $z$-axis of the new system,
  which results in the following vector:
  \[
    \mathbf{v}_{\mathrm{III}} =
    \begin{bmatrix}
      v_{\mathrm{Moon}} \cdot \cos(\kappa)\\
      v_{\mathrm{Moon}} \cdot \sin(\kappa) \cdot
        \cos(\theta_{\mathrm{umbra}} - \theta_{\mathrm{Moon}})\\
      -\,v_{\mathrm{Moon}} \cdot \sin(\kappa) \cdot
        \sin(\theta_{\mathrm{umbra}} - \theta_{\mathrm{Moon}})
    \end{bmatrix}
  \]
  \hfill(6.0)
  \[
    \mathbf{v}_{\mathrm{III}} =
    \begin{bmatrix}
      1085\ \mathrm{m/s}\\
      77.76\ \mathrm{m/s}\\
      -23.54\ \mathrm{m/s}
    \end{bmatrix}
  \]
  \hfill(1.0)

\item[(j)] The velocity of the umbra corresponds to the vector sum of the
  component caused by the rotation of the Earth and the component caused by
  the velocity of the Moon.

  Note that all points on the axis of the Moon's shadow cone have the same
  velocity, $\mathbf{v}_{\mathrm{III}}$ is also the velocity of the axis of
  the Moon's shadow cone at the center of the umbra. Since the umbra moves
  along the surface of the Earth, so the $z_{\mathrm{III}}$-component should
  be set to zero.
  \[
    \mathbf{v}_{\mathrm{umbra}} =
    \begin{bmatrix}
      v_{\mathrm{III},x} - v_{\mathrm{rot}}\\
      v_{\mathrm{III},y}\\
      0
    \end{bmatrix}
    =
    \begin{bmatrix}
      1085 - 463.7\\
      77.76\\
      0
    \end{bmatrix}
    =
    \begin{bmatrix}
      621.4\ \mathrm{m/s}\\
      77.59\ \mathrm{m/s}\\ % sic: 77.59 here, though the middle step and the modulus below use 77.76
      0
    \end{bmatrix}
  \]
  \hfill(4.0)

  The modulus of this vector is the following:
  \[
    v_{\mathrm{umbra}} = |\mathbf{v}_{\mathrm{umbra}}|
      = \sqrt{621.4^2 + 77.76^2} = 626.3\ \mathrm{m/s}
  \]
  \hfill(2.0)

\item[(k)] Considering the assumptions listed at the beginning of part II, it
  is possible to estimate the length of the totality by dividing the diameter
  of the umbra by its speed:
  \[
    \Delta t_{\mathrm{totality}}
      = \frac{2 \cdot r_{\mathrm{umbra}}}{v_{\mathrm{umbra}}}
      = \frac{2 \times 1.196 \times 10^{5}}{626.3}
  \]
  \hfill(2.0)
  \[
    \boxed{\Delta t_{\mathrm{totality}} = 6\,\mathrm{min}\,21.4\,\mathrm{s}}
  \]
  \hfill(1.0)
\end{description}

\solhead{3.29}{Galilean moons}
% source id: 2007_DA_D1
\ioaafig{2007_DA_D1-s1.pdf}
Names of the Galilean Moons

\begin{description}
\item[1.] Callisto \hfill(1.5 points)
\item[2.] Europa \hfill(1.5 points)
\item[3.] Io \hfill(1.5 points)
\item[4.] Ganymede \hfill(1.5 points)
% sic: the solution key awards 1.5 points each (6 total) although the exam paper titles the question "(4 points)"
\end{description}

\solhead{3.30}{Solar System objects}
% source id: 2007_DA_D3
\begin{description}
\item[3.1]
  \begin{center}
    \begin{tabular}{ll}
      Object   & Name    \\
      \hline
      Object A & Sun     \\
      Object B & Jupiter \\
      Object C & Venus   \\
      Object D & Mars    \\
    \end{tabular}
  \end{center}
  \hfill(0.5 points each)

\item[3.2] Object B \hfill(0.5 points)
  % sic: the marking scheme awards 0.5 points for 3.2 and 3.3 (7 total) although the exam paper says 1 point each (8 total)

\item[3.3] 3rd February \hfill(0.5 points)

\item[3.4]
\ioaafig{2007_DA_D3-s1.pdf}

  \begin{itemize}
  \item Correct position of the Earth relative to the Sun
    ($46^\circ$ before RA${}=0^h$) \hfill(1 point)
  \item Correct position of object in the orbit:\\
    Venus (C) ${}=0.5$ points\\
    Mars (D) ${}=0.5$ points\\
    Jupiter (B) ${}=0.5$ points
  \item Correct value of elongation angle:\\
    Venus (C) ${}=48.07^\circ$ \hfill(0.5 points)\\
    Mars (D) ${}=69.39^\circ$ \hfill(0.5 points)\\
    Jupiter (B) ${}=181.20^\circ$ % sic: elongation given as 181.20 degrees in the source
    \hfill(0.5 points)
  \end{itemize}
\end{description}

\solhead{3.31}{The Age of Meteorite}
% source id: 2008_DA_D3
\begin{description}
\item[1.] From $N(t) = N_0 \exp(-\lambda t)$ and $D(t) = N_0 - N(t)$, then we
  have
  \begin{align*}
    D(t) &= N(t)e^{\lambda t} - N(t) = N(t)(e^{\lambda t} - 1)\\
    \frac{D(t)}{N(t)} &= (e^{\lambda t} - 1)
      \ \text{ or }\ e^{\lambda t} = \frac{D(t)}{N(t)} + 1,\ \text{so}\\
    \lambda t &= \ln\left(\frac{D(t)}{N(t)} + 1\right) \ \text{ and}\\
    t &= \frac{1}{\lambda}\ln\left(\frac{D(t)}{N(t)} + 1\right)
  \end{align*}
  \hfill(10\%)

\item[2.] Setting $N(t) = \frac{1}{2}N_0$, we get:
  \[
    t_{1/2} = \ln 2/\lambda = 48.8 \times 10^{9}\ \mathrm{year}
            = 48.8\ \mathrm{Gyr}
  \]
  \hfill(5\%)

\item[3.] We know that $^{87}\mathrm{Rb} \rightarrow {}^{87}\mathrm{Sr}$, and
  taking the ratios $\frac{^{87}\mathrm{Sr}}{^{86}\mathrm{Sr}}$ as $y$
  (dependent variable) and $\frac{^{87}\mathrm{Rb}}{^{86}\mathrm{Sr}}$ as $x$
  (independent variable), we can use a simple linear regression equation:
  \[
    y = a + bx
  \]
  Taking into account the initial strontium, the corresponding equation is
  \[
    \left(\frac{^{87}\mathrm{Sr}}{^{86}\mathrm{Sr}}\right)
    = \left(\frac{^{87}\mathrm{Sr}}{^{86}\mathrm{Sr}}\right)_0
    + \left(e^{\lambda t} - 1\right)
      \left(\frac{^{87}\mathrm{Rb}}{^{86}\mathrm{Sr}}\right)
  \]
  where \textbf{the intercept}
  $a = \left(\frac{^{87}\mathrm{Sr}}{^{86}\mathrm{Sr}}\right)_0$;
  and \textbf{the slope} $b = \left(e^{\lambda t} - 1\right)$
  \hfill(20\%)

\item[4.] First, take the ratios:
  \begin{center}
    \begin{tabular}{ccccccc}
      Sample No & Meteorite type & $^{86}$Sr (ppm) & $^{87}$Rb (ppm)
        & $^{87}$Sr (ppm) & $^{87}$Rb/$^{86}$Sr & $^{87}$Sr/$^{86}$Sr\\
      \hline
      1  & A & 29.6 & 0.3   & 20.7 & 0.0101351 & 0.6993243\\
      2  & B & 58.7 & 68.5  & 44.7 & 1.1669506 & 0.7614991\\
      3  & B & 74.2 & 14.4  & 52.9 & 0.1940701 & 0.712938\\
      4  & A & 40.2 & 7.0   & 28.6 & 0.1741294 & 0.7114428\\
      5  & A & 19.7 & 0.4   & 13.8 & 0.0203046 & 0.7005076\\
      6  & B & 37.9 & 31.6  & 28.4 & 0.8337731 & 0.7493404\\
      7  & A & 33.4 & 4.0   & 23.6 & 0.1197605 & 0.7065868\\
      8  & B & 29.8 & 105.0 & 26.4 & 3.5234899 & 0.885906\\
      9  & A & 9.8  & 0.8   & 6.9  & 0.0816327 & 0.7040816\\
      10 & B & 18.5 & 44.0  & 15.4 & 2.3783784 & 0.8324324\\
    \end{tabular}
  \end{center}

  And we compile the data into Type A and B:

  \begin{center}
    \textbf{Type A}

    \medskip
    \begin{tabular}{cccccc}
      Sample No & $^{87}$Rb/$^{86}$Sr ($X$) & $^{87}$Sr/$^{86}$Sr ($Y$)
        & $XY$ & $X^2$ & $Y^2$\\
      \hline
      1 & 0.0101351 & 0.6993243 & 0.0070877 & 0.0001027 & 0.4890545\\
      4 & 0.1741294 & 0.7114428 & 0.1238831 & 0.030321  & 0.5061509\\
      5 & 0.0203046 & 0.7005076 & 0.0142235 & 0.0004123 & 0.4907109\\
      7 & 0.1197605 & 0.7065868 & 0.0846212 & 0.0143426 & 0.4992649\\
      9 & 0.0816327 & 0.7040816 & 0.0574761 & 0.0066639 & 0.4957309\\
      \hline
      Sum: & \textbf{0.4059623} & \textbf{3.5219431} & \textbf{0.2872916}
        & \textbf{0.0518425} & \textbf{2.4809121}\\
    \end{tabular}
  \end{center}

  \begin{align*}
    \lambda &= 1.42 \times 10^{-11}\\
    \bar{X} &= 0.0811925\\
    \bar{Y} &= 0.7043886\\
    SS_{xx} &= 0.0188814\\
    SS_{yy} &= 9.54601 \times 10^{-5}\\
    SS_{xy} &= 0.0013364\\
    b_A &= 0.0707786\\
    a_A &= 0.6986419
  \end{align*}
  where,
  \begin{align*}
    SS_{xx}&: \text{sum of square for } X
      = \sum_{i=1}^{n} X_i^2
        - \frac{1}{n}\left(\sum_{i=1}^{n} X_i\right)^{2}\\
    SS_{yy}&: \text{sum of square for } Y
      = \sum_{i=1}^{n} Y_i^2
        - \frac{1}{n}\left(\sum_{i=1}^{n} Y_i\right)^{2}\\
    SS_{xy}&: \text{sum of square for both } X \text{ and } Y
      = \sum_{i=1}^{n} X_i Y_i
        - \frac{1}{n}\sum_{i=1}^{n} X_i \sum_{i=1}^{n} Y_i
  \end{align*}
  and \textbf{the slope}
  \[
    b_A = \left(e^{\lambda t} - 1\right) = \frac{SS_{XY}}{SS_{XX}}
    = \frac{\displaystyle\sum_{i=1}^{5} X_i Y_i
        - \frac{1}{5}\sum_{i=1}^{5} X_i \sum_{i=1}^{5} Y_i}
      {\displaystyle\sum_{i=1}^{5} X_i^2
        - \frac{1}{5}\left(\sum_{i=1}^{5} X_i\right)^{2}}
    = \frac{0.2872916 - 0.2859552}{0.0518425 - 0.0329611}
    = 0.0707786 \qquad (*)
  \]
  and \textbf{the intercept}
  \begin{align*}
    a_A &= \left(\frac{^{87}\mathrm{Sr}}{^{86}\mathrm{Sr}}\right)_0
      = \bar{Y} - b\bar{X}
      = \overline{\left(\frac{^{87}\mathrm{Sr}}{^{86}\mathrm{Sr}}\right)}
        - b\,\overline{\left(\frac{^{87}\mathrm{Rb}}{^{86}\mathrm{Sr}}\right)}\\
      &= 0.7043886 - 0.0707786 \times 0.0811925 = 0.6986419
  \end{align*}

  The scatter plot and linear regression line for Type A is:

\ioaafig{2008_DA_D3-s1.pdf}
  \hfill(20\%)

  \begin{center}
    \textbf{Type B}

    \medskip
    \begin{tabular}{cccccc}
      Sample No & $^{87}$Rb/$^{86}$Sr ($X$) & $^{87}$Sr/$^{86}$Sr ($Y$)
        & $XY$ & $X^2$ & $Y^2$\\
      \hline
      2  & 1.1669506 & 0.7614991 & 0.8886318 & 1.3617737  & 0.5798809\\
      3  & 0.1940701 & 0.7129380 & 0.1383599 & 0.0376632  & 0.5082806\\
      6  & 0.8337731 & 0.7493404 & 0.6247799 & 0.6951776  & 0.561511\\
      8  & 3.5234899 & 0.8859060 & 3.1214808 & 12.4149811 & 0.7848294\\
      10 & 2.3783784 & 0.8324324 & 1.9798392 & 5.6566838  & 0.6929437\\
      \hline
      Sum: & \textbf{8.0966621} & \textbf{3.9421159} & \textbf{6.7530916}
        & \textbf{20.1662794} & \textbf{3.1274456}\\
    \end{tabular}
  \end{center}

  \begin{align*}
    \lambda &= 1.42 \times 10^{-11}\\
    \bar{X} &= 1.6193324\\
    \bar{Y} &= 0.7884232\\
    SS_{xx} &= 7.0550920\\
    SS_{yy} &= 0.019390046\\
    SS_{xy} &= 0.3694955\\
    b_B &= 0.0523729\\
    a_B &= 0.7036141
  \end{align*}
  and \textbf{the slope} $b_B = \left(e^{\lambda t} - 1\right)$
  \[
    = \frac{SS_{XY}}{SS_{XX}}
    = \frac{\displaystyle\sum_{i=1}^{5} X_i Y_i
        - \frac{1}{5}\sum_{i=1}^{5} X_i \sum_{i=1}^{5} Y_i}
      {\displaystyle\sum_{i=1}^{5} X_i^2
        - \frac{1}{5}\left(\sum_{i=1}^{5} X_i\right)^{2}}
    = \frac{6.7530916 - 6.3835961}{20.1662794 - 13.1111874}
    = 0.0523729
  \]
  and \textbf{the intercept}
  \begin{align*}
    a_B &= \left(\frac{^{87}\mathrm{Sr}}{^{86}\mathrm{Sr}}\right)_0
      = \bar{Y} - b_B\bar{X}
      = \overline{\left(\frac{^{87}\mathrm{Sr}}{^{86}\mathrm{Sr}}\right)}
        - b_B\,\overline{\left(\frac{^{87}\mathrm{Rb}}{^{86}\mathrm{Sr}}\right)}
      = 0.7884232 - 0.0523729 \times 1.6193324\\
      &= 0.7036141
  \end{align*}

  The scatter plot and linear regression line for Type B is:

\ioaafig{2008_DA_D3-s2.pdf}
  \hfill(20\%)

\item[5.] Since we have 2 types of meteorites, then we have to calculate their
  ages separately, and $\lambda = 1.42 \times 10^{-11}$. As
  $b = \left(e^{\lambda t} - 1\right)$, then $e^{\lambda t} = b + 1$.\\
  After we get $a$ and $b$ then
  \[
    t = \frac{1}{\lambda}\ln(b + 1)
  \]
  which is $4.81592 \times 10^{9}$ year.

  From $(*)$, we obtain
  \begin{align*}
    s_b &= 0.0040112\\
    s_a &= 0.0009133\\
    e_t &= 2.819138 \times 10^{8}
  \end{align*}
  sigma-$b$:
  \[
    s_b = \sqrt{\frac{SS_{YY} - \dfrac{(SS_{XY})^2}{SS_{XX}}}
                     {(n-2)\,SS_{XX}}}
        = \sqrt{\frac{0.0000955 - \dfrac{(0.0013364)^2}{0.0188814}}
                     {3 \times 0.0188814}}
        = 0.0040112
  \]
  error on $t$:
  \[
    e_t = \frac{1}{\lambda}\ln(s_b + 1)
        = \frac{\ln(0.0040112 + 1)}{1.42 \times 10^{-11}}
        = 2.819138 \times 10^{8}
  \]
  sigma-$a$ (error on initial):
  \[
    s_a = \sqrt{\frac{SS_{YY} - \dfrac{(SS_{XY})^2}{SS_{XX}}}
                     {(n-2)\,SS_{XX}} \times \sum_{i=1}^{n} X_i^2}
        = \sqrt{\frac{0.0000955 - \dfrac{(0.0013364)^2}{0.0188814}}
                     {3 \times 0.0188814} \times 0.0518425}
        = 0.0009133
  \]

  \textbf{Type B}

  Similar as in Type A, we obtain for Type B, in which
  $t = 3.595 \times 10^{9}$ year.
  \begin{align*}
    s_b &= 0.0013487\\ % sic: 0.0013487 listed, though the formula below gives 0.0013479
    s_a &= 0.0060566\\ % sic: 0.0060566 listed, though the formula below gives 0.0060532
    e_t &= 94915372.7530205
  \end{align*}
  sigma-$b$:
  \[
    s_b = \sqrt{\frac{SS_{YY} - \dfrac{(SS_{XY})^2}{SS_{XX}}}
                     {(n-2)\,SS_{XX}}}
        = \sqrt{\frac{0.0193900 - \dfrac{0.3694955^2}{7.0550920}}
                     {3 \times 7.0550920}}
        = 0.0013479
  \]
  error on $t$:
  \[
    e_t = \frac{1}{\lambda}\ln(s_b + 1)
        = \frac{\ln(0.0013479 + 1)}{1.42 \times 10^{-11}}
        = 9.48586 \times 10^{7}
  \]
  sigma-$a$ (error on initial):
  \[
    s_a = \sqrt{\frac{SS_{YY} - \dfrac{(SS_{XY})^2}{SS_{XX}}}
                     {(n-2)\,SS_{XX}} \times \sum_{i=1}^{n} X_i^2}
        = \sqrt{\frac{0.0193900 - \dfrac{0.3694955^2}{7.0550920}}
                     {3 \times 7.0550920} \times 20.1662794}
        = 0.0060532
  \]

  The Age of type A ${}= (4.816 \pm 0.282) \times 10^{9}$ year
  ${}= 4.816 \pm 0.282$ Gyr

  The Age of type B ${}= (3.595 \pm 0.095) \times 10^{9}$ year
  ${}= 3.595 \pm 0.435$ Gyr % sic: 0.095 and 0.435 both appear in the source

  Hence, the type A is much older than the type B.
  \hfill(20\%)

\item[6.] The initial value of type A
  $\left(^{87}\mathrm{Sr}/^{86}\mathrm{Sr}\right)_0
    = 0.6986419 \pm 0.0009133$\\
  The initial value of type B
  $\left(^{87}\mathrm{Sr}/^{86}\mathrm{Sr}\right)_0
    = 0.7036141 \pm 0.0060532$
  \hfill(5\%)
\end{description}

\solhead{3.32}{Venus}
% source id: 2009_DA_D2
\begin{description}
\item[a)]
The $\angle SVE$ angle should be calculated from the phase of Venus. Figure 2.1 shows that
projected area of Venus disk which is illuminated by the Sun is
\[
\frac{\pi r^2}{2} + \frac{\pi r r'}{2}
\]
where
\[
r' = r\cos(\angle SVE)
\]
Then,
\[
Phase = \Bigl(\frac{\dfrac{\pi r^2}{2} + \dfrac{\pi r^2\cos(\angle SVE)}{2}}{\pi r^2}\Bigr)
\times 100
= \frac{100}{2}\,(1 + \cos(\angle SVE)) = 100\cos^2\Bigl(\frac{\angle SVE}{2}\Bigr)
\]
The angle $\angle SVE$ is calculated and written in table 2.2, column 2.
\ioaafig{2009_DA_D2-s1.pdf}

\item[b)]
As in figure 2.1, in $SEV$ triangle we have,
\begin{align*}
\frac{r_e}{\sin(\angle SVE)} &= \frac{r_v}{\sin(\angle SEV)}\\
r_v &= r_e\,\frac{\sin(\angle SEV)}{\sin(\angle SVE)}
\end{align*}
where $r_e$ and $\angle SEV$ (elongation) is given in table 2.1 then, $r_v$ for all observing
dates is calculated and written in table 2.2 column 3.

\item[c)]
\ioaafig{2009_DA_D2-s2.pdf}

\item[d)]
According to the obtained values written in table 2.2 column 3,
\begin{align*}
r_v^{max} &= 0.728\ AU\\
r_v^{min} &= 0.718\ AU
\end{align*}

\item[e)]
Semi-major axis is
\[
a = \frac{(r_v^{max} + r_v^{min})}{2} = 0.723\ AU
\]

\item[f)]
Eccentricity could be calculated from both of aphelion and perihelion distances as
\[
e = \frac{r_v^{max} - r_v^{min}}{2a} = 6.92\times10^{-3}
\]
\end{description}

\begin{center}
Table 2.2

\begin{tabular}{|c|c|c|}
\hline
Column 1 & Column 2 & Column 3\\
\hline
Date & SVE ($^\circ$) & Sun - Venus Distance (AU)\\
\hline
2008-Sep-20 & 39.83 & 0.726\\
\hline
2008-Oct-10 & 47.16 & 0.728\\
\hline
2008-Oct-20 & 50.80 & 0.728\\
\hline
2008-Oct-30 & 54.55 & 0.728\\
\hline
2008-Nov-09 & 58.26 & 0.728\\
\hline
2008-Nov-19 & 62.10 & 0.728\\
\hline
2008-Nov-29 & 66.17 & 0.727\\
\hline
2008-Dec-19 & 74.81 & 0.725\\
\hline
2008-Dec-29 & 79.63 & 0.723\\
\hline
2009-Jan-18 & 90.57 & 0.721\\
\hline
2009-Feb-07 & 104.83 & 0.719\\
\hline
2009-Feb-17 & 114.08 & 0.718\\
\hline
2009-Feb-27 & 125.59 & 0.719\\
\hline
2009-Mar-19 & 157.52 & 0.721\\
\hline
\end{tabular}
\end{center}

\textbf{Note:} reported numbers in table 2 are not acceptable if they are out of 0.75 and
1.25 times of designer answer.
% The standalone "Venus Problem Marking Scheme" table (p.77) is omitted.

\solhead{3.33}{Asteroid photometry}
% source id: 2012_DA_D1
% Official solution (2012_da_sol.pdf, "Data analysis", August 10, 2012, Question 1).
% The published solution consists of a point distribution ("Pontuation", omitted here as a
% standalone marking scheme) and the answer list below; all quantitative results are
% presented in Figures 1-7.

Answers:
\begin{itemize}
\item Items 1,2 \textbf{See Fig. 1--3.}
\item Item 3: \textbf{Night B}
\item Items 4,5,6 \textbf{See Fig. 4--6.}
\item Items 7,8 \textbf{See Fig. 7.}
\item Item 9 \textbf{Yes}
\end{itemize}

\ioaafig{2012_DA_D1-s1.pdf}
\ioaafig{2012_DA_D1-s2.pdf}
\ioaafig{2012_DA_D1-s3.pdf}
\ioaafig{2012_DA_D1-s4.pdf}
\ioaafig{2012_DA_D1-s5.pdf}
\ioaafig{2012_DA_D1-s6.pdf}
\ioaafig{2012_DA_D1-s7.pdf}

\solhead{3.34}{Thermodynamic test}
% source id: 2014_DA_D2
% Official solution ("Problem 3 A. Marking scheme -- Thermodynamic test, 25 points",

\begin{description}
\item[a)]
The pressure and the density of the atmosphere increase with the decrease of the altitude.
The pressure at the surface of the planet is $p_{\max}$.

In figure 1 is represented a thin atmospheric layer with thickness $\Delta h$ near the surface
of the planet P$_1$ where the atmosphere density can be considered $\rho_{\max}$. Inside this
layer the pressure and the temperature are variable, reaching the values $p_{\max}$ and $T_0$
near the surface.
\ioaafig{2014_DA_D2-s1.pdf}

The following relations can be written:
\begin{align*}
p_{\max} &= p + \Delta p;\\
\Delta p &= \rho_{\max}g_0\Delta h,
\end{align*}
$\Delta p$ -- the variation of the pressure due to the altitude variation $\Delta h$, $g_0$
the gravitational acceleration near the surface of planet P$_1$;
\[
\Delta h = \frac{\Delta p}{\rho_{\max}g_0}.
\]
The superior layer where the density is constant $\rho$ is represented too. At its base the
pressure and temperature hav the values $p$ and respectively $T$. % sic: "hav"

For a gas volume inside the near surface layer of the atmosphere, near the surface of the
planet for $\nu$ mol of de gas according with Mendeleev--Klapeyron law % sic: "of de gas"
\begin{align*}
p_{\max}V &= \nu RT_0 = \frac{m}{\mu}RT_0;\\
p_{\max}\mu &= \frac{m}{V}RT_0 = \rho_{\max}RT_0;\\
\rho_{\max} &= \frac{p_{\max}\mu}{RT_0}.
\end{align*}

Just before touching the surface of the planet P$_1$ the radio-probe moves uniformly, so the
distance $\Delta h$ is spend with speed $\vec{\mathrm{v}}$, in time $\Delta t$: % sic: "is spend"
\begin{align*}
\Delta h &= \mathrm{v}\,\Delta t;\\
\Delta h &= \frac{\Delta p}{\rho_{\max}g_0};\quad
\rho_{\max} = \frac{p_{\max}\mu}{RT_0};\\
\frac{\Delta p}{\rho_{\max}g_0} &= \mathrm{v}\,\Delta t;\quad
\mathrm{v} = \frac{\Delta p}{\rho_{\max}g_0\Delta t};\quad
\mathrm{v} = \frac{\Delta p}{\dfrac{p_{\max}\mu}{RT_0}g_0\Delta t};\\
\mathrm{v} &= \frac{RT_0\Delta p}{g_0\mu p_{\max}\Delta t},\quad(1)
\end{align*}
Where accordingly to fig.~2 in the point F $(\tau;\,p_{\max})$, the trigonometric tangent of
$\alpha$ is:
\[
\tan\alpha = \frac{\Delta p}{\Delta t}.
\]
\ioaafig{2014_DA_D2-s2.pdf}

The rapport: % sic: "rapport"
\[
\frac{\Delta p}{p_{\max}\Delta t},
\]
Can be determined in the graph in the point F:
\[
\mathrm{v} = \frac{\Delta p}{p_{\max}\Delta t}\,\frac{RT_0}{g_0\mu};
\]
\[
\frac{\Delta p}{p_{\max}\Delta t}
= \frac{15\,\mathrm{u.c.}}{55\,\mathrm{u.c.}\,250\,\mathrm{s}}
= \frac{3}{2750}\,\frac{1}{\mathrm{s}}
\qquad
\frac{45\,\mathrm{u.c.}}{55\,\mathrm{u.c.}\,750\,\mathrm{s}}
= \frac{3}{2750}\,\frac{1}{\mathrm{s}};
\]
\begin{align*}
\mathrm{v} &= \frac{3}{2750}\,\frac{1}{\mathrm{s}}\,
\frac{8.3\,\dfrac{\mathrm{J}}{\mathrm{mol\,K}}\;700\,\mathrm{K}}
     {10\,\dfrac{\mathrm{m}}{\mathrm{s}^2}\;44\,\dfrac{10^{-3}\,\mathrm{kg}}{\mathrm{mol}}}
= 14.4\,\frac{\mathrm{m}}{\mathrm{s}};\\
h_0 &= \mathrm{v}\tau;\quad \tau = 3.750\ \mathrm{s};\\
h_0 &= 54.000\,\mathrm{m} = 54\,\mathrm{km}.
\end{align*}
% Dots in 3.750 s and 54.000 m are thousands separators, as in the original.

\item[b)]
Descending from the initial altitude $h_0 = 54\,\mathrm{km}$, with speed
$\mathrm{v} = 14.4\,\frac{\mathrm{m}}{\mathrm{s}}$, the radio-probe, after 1000 s arrives
at the altitude $h = 39.6\,\mathrm{km}$
\[
t = \frac{h_0 - h}{\mathrm{v}} = \frac{14400\,\mathrm{m}}{14.4\,\frac{\mathrm{m}}{\mathrm{s}}}
= 1000\,\mathrm{s}.
\]
On Graph at altitude $h = 39.6\,\mathrm{km}$, the atmospheric pressure is
$p_h = 4\,\mathrm{u.c.}$, -- see point C on grah. % sic: "grah"
\ioaafig{2014_DA_D2-s3.pdf}

In this point
\begin{align*}
\tan\beta &= \left(\frac{\Delta p}{\Delta t}\right)_{\!C};\\
\left(\frac{\Delta p}{p\Delta t}\right)_{\!C}
&= \frac{1}{p_h}\left(\frac{\Delta p}{\Delta t}\right)_{\!C}
= \frac{1}{4\,\mathrm{u.c.}}\,\frac{3\,\mathrm{u.c.}}{500\,\mathrm{s}}
= \frac{3}{2000}\,\frac{1}{\mathrm{s}};\\
\left(\frac{p\Delta t}{\Delta p}\right)_{\!C}
&= p_h\left(\frac{\Delta t}{\Delta p}\right)_{\!C}
= 4\,\mathrm{u.c.}\,\frac{500\,\mathrm{s}}{3\,\mathrm{u.c.}}
= \frac{2000}{3}\,\mathrm{s};\\
\mathrm{v} &= \left(\frac{\Delta p}{p\Delta t}\right)_{\!C}\frac{RT_h}{g_0\mu};\\
\mathrm{v} &= \frac{1}{p_h}\left(\frac{\Delta p}{\Delta t}\right)_{\!C}\frac{RT_h}{g_0\mu};\\
T_h &= \frac{\mathrm{v}\,g_0\mu}
{\dfrac{1}{p_h}\left(\dfrac{\Delta p}{\Delta t}\right)_{\!C}R};\quad
T_h = p_h\left(\frac{\Delta t}{\Delta p}\right)_{\!C}\frac{\mathrm{v}\,g_0\mu}{R};\\
T_h &= \frac{2000}{3}\,\mathrm{s}\,
\frac{14.4\,\dfrac{\mathrm{m}}{\mathrm{s}}\;10\,\dfrac{\mathrm{m}}{\mathrm{s}^2}\;
      44\,\dfrac{10^{-3}\,\mathrm{kg}}{\mathrm{mol}}}
     {8.3\,\dfrac{\mathrm{J}}{\mathrm{mol\,K}}};\\
T_h &= 508.9\,\mathrm{K}.
\end{align*}

\item[c)]
By using the same algorithm, corresponding to the points F, D, C, B \c{s}i A in graph on
figure 4: % sic: Romanian "şi" (and) left in the original
\ioaafig{2014_DA_D2-s4.pdf}
\begin{align*}
\mathrm{v} &= \frac{\Delta p}{p_{\max}\Delta t}\,\frac{RT_0}{g_0\mu};\\
\mathrm{v} &= \frac{10\,\mathrm{u.c.}}{50\,\mathrm{u.c.}\,500\,\mathrm{s}}\,
\frac{8.3\,\dfrac{\mathrm{J}}{\mathrm{mol\,K}}\;750\,\mathrm{K}}
     {8\,\dfrac{\mathrm{m}}{\mathrm{s}^2}\;44\,\dfrac{10^{-3}\,\mathrm{kg}}{\mathrm{mol}}}
\approx 7\,\frac{\mathrm{m}}{\mathrm{s}};\\
h_0 &= \mathrm{v}\tau = 7\,\frac{\mathrm{m}}{\mathrm{s}}\,2000\,\mathrm{s}
= 14.000\,\mathrm{m} = 14\,\mathrm{km};\\
h_F &= 0;
\end{align*}
\[
\mathrm{F}(p = 50\,\mathrm{u.c.};\ h = 0\,\mathrm{km};\ T = 750\,\mathrm{K});
\]
\begin{align*}
T_D &= p_D\left(\frac{\Delta t}{\Delta p}\right)_{\!D}\frac{\mathrm{v}\,g_0\mu}{R}
= 593.7\,\mathrm{K};\\
h_D &= h_0 - 7\,\frac{\mathrm{m}}{\mathrm{s}}\,1500\,\mathrm{s}
= 14\,\mathrm{km} - 10.500\,\mathrm{m} = 3.5\,\mathrm{km};
\end{align*}
\[
\mathrm{D}(p = 40\,\mathrm{u.c.};\ h = 3.5\,\mathrm{km};\ T = 593.7\,\mathrm{K});
\]
\begin{align*}
T_C &= p_C\left(\frac{\Delta t}{\Delta p}\right)_{\!C}\frac{\mathrm{v}\,g_0\mu}{R}
= 445.3\,\mathrm{K};\\
h_C &= h_0 - 7\,\frac{\mathrm{m}}{\mathrm{s}}\,1000\,\mathrm{s}
= 14\,\mathrm{km} - 7.000\,\mathrm{m} = 7\,\mathrm{km};
\end{align*}
\[
\mathrm{C}(p = 30\,\mathrm{u.c.};\ h = 7\,\mathrm{km};\ T = 445.3\,\mathrm{K});
\]
\begin{align*}
T_B &= p_B\left(\frac{\Delta t}{\Delta p}\right)_{\!B}\frac{\mathrm{v}\,g_0\mu}{R}
= 296.8\,\mathrm{K};\\
h_B &= h_0 - 7\,\frac{\mathrm{m}}{\mathrm{s}}\,500\,\mathrm{s}
= 14\,\mathrm{km} - 3.500\,\mathrm{m} = 10.5\,\mathrm{km};
\end{align*}
\[
\mathrm{B}(p = 20\,\mathrm{u.c.};\ h = 10.5\,\mathrm{km};\ T = 296.8\,\mathrm{K});
\]
\begin{align*}
T_A &= p_A\left(\frac{\Delta t}{\Delta p}\right)_{\!A}\frac{\mathrm{v}\,g_0\mu}{R}
= 148.4\,\mathrm{K};\\
h_A &= 14\,\mathrm{km};
\end{align*}
\[
\mathrm{A}(p = 10\,\mathrm{u.c.};\ h = 14\,\mathrm{km};\ T = 148.4\,\mathrm{K}).
\]

The data are in table below which allows to plot de graph $p = f(h)$ and $T = f(h)$ for
planet P$_2$ figure 5 and figure 6. % sic: "plot de graph"
\begin{center}
\begin{tabular}{|c|c|c|c|c|c|}
\hline
$t$ & 0 & 500 s & 1000 s & 1500 s & 2000 s\\
\hline
$p$ & 10 u.c. & 20 u.c. & 30 u.c. & 40 u.c. & 50 u.c.\\
\hline
$h$ & 14 km & 10,5 km & 7 km & 3,5 km & 0\\
\hline
$T$ & 148,4 K & 296,8 K & 445,3 K & 593,3 K & 750 K\\
\hline
\end{tabular}
\end{center}
% sic: the table prints 593,3 K although the computation above gives 593,7 K.
\ioaafig{2014_DA_D2-s5.pdf}
\ioaafig{2014_DA_D2-s6.pdf}
\end{description}

\solhead{3.35}{Distance to the Moon}
% source id: 2016_DA_D2
\begin{description}
\item[(D2.1)]
From the table we see that the angular size of Moon is smallest close to the New Moon day.
Thus, the answer is \fbox{New Moon}. \hfill(3.0)

\emph{Justification is NOT necessary for full credit.}

\item[(D2.2)]
As there is an eclipse happening in this month, the lunar nodes are close to Full Moon day
and New Moon day. Next we notice that lowest declination of Moon is just $18^\circ$. This
means that after the New Moon day, the orbit of Moon is above the ecliptic. In other words,
the ascending node is near the \fbox{New Moon}. \hfill(4.0)

\emph{Justification is NOT necessary for full credit.}

\item[(D2.3)]
The largest angular size of the Moon in the ephemerides is $2008.3''$ and the smallest
angular size is $1763.7''$. The distance is inversely proportional to the angular size. Hence
ratio of distance at perigee to the distance at apogee is:
\begin{align*}
\mathrm{Ratio} = \frac{r_{\mathrm{perigee}}}{r_{\mathrm{apogee}}}
 &= \frac{a_0}{1+e}\times\frac{1-e}{a_0}
 & &\hspace{-2em}(2.0)\\
 % the two a_0 factors are struck through (cancelled) in the original
\therefore \frac{1-e}{1+e} &= \frac{1763.7}{2008.3} = 0.87821\\
\therefore e &= \frac{1 - 0.87821}{1 + 0.87821} = 0.064846\\
&\fbox{$e \simeq 0.065$}
 & &\hspace{-2em}(2.0)
\end{align*}
\emph{Rounding off is done to account for the fact that our data are not continuous, hence
exact angular sizes at perigee and apogee are not known. A non-rounded answer will also
receive full credit.}

\item[(D2.4)]
The following construction is shown. \hfill(5.0)
\ioaafig{2016_DA_D2-s1.pdf}

Only two chords are necessary to determine the centre.

\emph{The credit is divided in two parts for the drawing:}
\begin{itemize}
\item \emph{Realisation that centre of umbra circle needs to be determined to find
$\theta_{\mathrm{umbra}}$: 1.5}
\item \emph{Accurate determination of centre of umbra circle by geometric construction: 2.5\\
Determination of centre of hand-drawn circle: maximum 1.5}
\item \emph{Measuring diameters of umbra and Moon: 0.5 each}
\end{itemize}

By estimating approximate centre of the shadow in the image, we find out,
\begin{align*}
\frac{\theta_{\mathrm{umbra}}}{\theta_{\mathrm{Moon}}}
 &= \frac{d_{\mathrm{umbra}}}{d_{\mathrm{Moon}}} & &\hspace{-2em}(2.0)\\
 &= \frac{9.1}{3.3} = 2.76\\
\therefore \theta_{\mathrm{umbra}} &\simeq 2.76\,\theta_{\mathrm{Moon}}
 & &\hspace{-2em}(1.0)
\end{align*}
\emph{Acceptable range: $\pm0.10$.}

\item[(D2.5)]
The following diagram needs to be drawn.
\ioaafig{2016_DA_D2-s2.pdf}

Angular size of umbra is $\theta_{\mathrm{umbra}} = 2\angle BOQ$\\
Angular size of penumbra is $\theta_{\mathrm{penumbra}} = 2\angle AOQ$

We have
\begin{align*}
QA &= QP + PA\\
 &= OE + PE\tan\theta_1\\
 &\approx R_\oplus + d_{\mathrm{Moon}}\theta_1 & &\text{since } PA \ll PE\\
 &\approx R_\oplus + d_{\mathrm{Moon}}\frac{\theta_{\mathrm{Sun}}}{2}
   & &\text{since } \theta_1 \approx \theta_2 \approx \theta_{\mathrm{Sun}}/2
   \quad(2.0)
\end{align*}
and
\begin{align*}
QB &= QP - PB\\
 &= OE - PE\tan\theta_2\\
 &\approx R_\oplus - d_{\mathrm{Moon}}\theta_2\\
 &\approx R_\oplus - d_{\mathrm{Moon}}\frac{\theta_{\mathrm{Sun}}}{2}
   & &\hspace{-2em}(2.0)
\end{align*}
\begin{align*}
\therefore \theta_{\mathrm{penumbra}} = 2\angle AOQ
 = 2\tan^{-1}\left(\frac{QA}{OQ}\right) &\approx 2\frac{QA}{OQ}
   \qquad\text{since } QA \ll OQ\\
 = 2\frac{R_\oplus + d_{\mathrm{Moon}}\dfrac{\theta_{\mathrm{Sun}}}{2}}{d_{\mathrm{Moon}}}
 &= \frac{2R_\oplus}{d_{\mathrm{Moon}}} + \theta_{\mathrm{Sun}}
   & &\hspace{-2em}(1.0)
\end{align*}
and
\begin{align*}
\theta_{\mathrm{umbra}} = 2\angle BOQ
 = 2\tan^{-1}\left(\frac{QB}{OQ}\right) &\approx 2\frac{QB}{OQ}\\
 = 2\frac{R_\oplus - d_{\mathrm{Moon}}\dfrac{\theta_{\mathrm{Sun}}}{2}}{d_{\mathrm{Moon}}}
 &= \frac{2R_\oplus}{d_{\mathrm{Moon}}} - \theta_{\mathrm{Sun}}
   & &\hspace{-2em}(1.0)
\end{align*}
Subtracting,
\[
\theta_{\mathrm{penumbra}} - \theta_{\mathrm{umbra}} = 2\theta_{\mathrm{Sun}}
\Rightarrow \theta_{\mathrm{penumbra}} = \theta_{\mathrm{umbra}} + 2\theta_{\mathrm{Sun}}
\]
\hfill(1.0)

We have,
\[
\theta_{\mathrm{umbra}} = 2.76\,\theta_{\mathrm{Moon}}
\qquad\text{and } \theta_{\mathrm{Sun}} = 1915.0''
\]
From the given data, $\theta_{\mathrm{Moon}} = 2008.3''$. \hfill(1.0)

Therefore,
\begin{align*}
\theta_{\mathrm{penumbra}} &= 2.76\,\theta_{\mathrm{Moon}}
 + 2\frac{1915.0}{2008.3}\theta_{\mathrm{Moon}}\\
&\fbox{$\theta_{\mathrm{penumbra}} = 4.67\,\theta_{\mathrm{Moon}}$}
 & &\hspace{-2em}(1.0)
\end{align*}
\emph{Acceptable range: $4.57\,\theta_{\mathrm{Moon}}$ to $4.77\,\theta_{\mathrm{Moon}}$.}

\underline{Alternative solution}:

In the figure below, rays $HEA$ and $IFC$ are coming from one edge of solar disk and rays
$HFD$ and $GEB$ are coming from the opposite edge. The observer ($O$) is assumed to be at
the centre of the Earth. The Moon travels along the path $ABCD$ during the course of
eclipse.
\ioaafig{2016_DA_D2-s3.pdf}

From figure,
\[
\angle AEB = \angle GEH = \angle HFI = \angle DFC = \angle EJF = \theta_{\mathrm{Sun}}
\]
\hfill(3.0)
\begin{align*}
\theta_{\mathrm{umbra}} &= \angle BOC = 2.76\,\theta_{\mathrm{Moon}}
 & &\hspace{-2em}(1.0)\\
\theta_{\mathrm{penumbra}} &= \angle AOD & &\hspace{-2em}(1.0)\\
\angle AOD &= \angle AOB + \angle BOC + \angle COD\\
\angle AOB &= \angle AEB & &\hspace{-2em}(1.0)\\
\angle COD &= \angle CFD & &\hspace{-2em}(1.0)\\
\theta_{\mathrm{penumbra}} &\simeq \angle AEB + \theta_{\mathrm{umbra}} + \angle CFD\\
 &= 2\theta_{\mathrm{Sun}} + 2.76\,\theta_{\mathrm{Moon}}
 = 2\times1915.0'' + 2.76\times2008.3''\\
\theta_{\mathrm{penumbra}} &= 9372.9'' = 4.67\,\theta_{\mathrm{Moon}}
 & &\hspace{-2em}(2.0)
\end{align*}

\item[(D2.6)]
From the Moon,
\[
\theta_{\mathrm{Earth}} = \frac{2R_\oplus}{d_{\mathrm{Moon}}}
\]
\hfill(1.0)

From part (D2.5),
\[
\theta_{\mathrm{umbra}} + \theta_{\mathrm{penumbra}} = 2\theta_{\mathrm{Earth}}
\]
\hfill(2.0)
\begin{align*}
\therefore \theta_{\mathrm{Earth}}
 &= \frac{\theta_{\mathrm{umbra}} + \theta_{\mathrm{penumbra}}}{2}
 = \frac{2.76 + 4.67}{2}\theta_{\mathrm{Moon}} = 3.72\,\theta_{\mathrm{Moon}}\\
&\fbox{$\theta_{\mathrm{Moon}} = {} = 0.269\,\theta_{\mathrm{Earth}}$}
 & &\hspace{-2em}(2.0)
\end{align*}
% sic: the boxed answer prints a doubled equals sign.

\underline{Alternative solution}:

Let us say that the Moon is at position of $B$. Thus, angular size of Earth as seen from
this position will be, (see figure in the previous part)
\begin{align*}
\theta_{\mathrm{Earth}} &= \angle EBF = \angle BFD & &\hspace{-2em}(2.0)\\
 &= \angle BFC + \angle CFD\\
 &\simeq \theta_{\mathrm{umbra}} + \theta_{\mathrm{Sun}} & &\hspace{-2em}(1.0)
\end{align*}
The angular size of the Full Moon on 28 September as seen in the table is $2008.3''$.
\begin{align*}
\theta_{\mathrm{Earth}} &= 2.76\times2008.3'' + 1915.0'' = 7453.0''\\
&\fbox{$\theta_{\mathrm{Moon}} = 0.269\,\theta_{\mathrm{Earth}}
 = \dfrac{\theta_{\mathrm{Earth}}}{3.72}$}
 & &\hspace{-2em}(2.0)
\end{align*}

\item[(D2.7)]
Thus, the radius of Moon will be,
\begin{align*}
R_{\mathrm{Moon}} &= \frac{R_\oplus}{3.72} & &\hspace{-2em}(1.0)\\
R_{\mathrm{Moon}} &= \frac{6371}{3.72}\\
&\fbox{$R_{\mathrm{Moon}} \simeq 1713\,\mathrm{km}$} & &\hspace{-2em}(2.0)
\end{align*}
\emph{Acceptable range: $\pm20\,\mathrm{km}$.}

\item[(D2.8)]
The shortest and longest distances will be,
\begin{align*}
r_{\mathrm{perigee}} &= \frac{2\times1713\times206265}{2008.3}\\
&\fbox{$r_{\mathrm{perigee}} = 3.52\times10^5\,\mathrm{km}$} & &\hspace{-2em}(2.0)\\
r_{\mathrm{apogee}} &= \frac{2\times1713\times206265}{1763.7}\\
&\fbox{$r_{\mathrm{apogee}} = 4.01\times10^5\,\mathrm{km}$} & &\hspace{-2em}(2.0)
\end{align*}

\item[(D2.9)]
\ioaafig{2016_DA_D2-s4.pdf}
On September 10, phase of Moon is 0.097 and elongation of Moon is $36.2^\circ$. Angular
size of the Moon on this day is $1792.0''$. Therefore, distance to Moon (from Earth) on
September 10 is
\begin{align*}
d_{\mathrm{Moon,10}} &= \frac{2\times1713\times206265}{1792.0}\\
 &= 3.94\times10^5\,\mathrm{km} & &\hspace{-2em}(2.0)
\end{align*}
Let
\begin{align*}
\angle EMS &= \varphi_M\\
\angle ESM &= \varphi_S\\
\angle SEM &= \varphi_E\\
\therefore \varphi_E &= 36.2^\circ & &\hspace{-2em}(1.0)\\
\mathrm{phase} &= \frac{1 + \cos\varphi_M}{2} & &\hspace{-2em}(2.0)\\
\therefore \varphi_M &= \cos^{-1}\left(2\times\mathrm{phase} - 1\right)\\
 &= \cos^{-1}\left(2\times0.097 - 1\right) = \cos^{-1}(-0.806)\\
 &= 143.71^\circ & &\hspace{-2em}(2.0)\\
\varphi_S &= 180^\circ - \varphi_E - \varphi_M\\
 &= 180^\circ - 36.2^\circ - 143.71^\circ\\
 &= 0.09^\circ & &\hspace{-2em}(1.0)
\end{align*}
Now using sine rule,
\begin{align*}
\frac{d_{\mathrm{Sun}}}{d_{\mathrm{Moon,10}}} &= \frac{\sin\varphi_M}{\sin\varphi_S}
 & &\hspace{-2em}(2.0)\\
\therefore d_{\mathrm{Sun}} &= 3.94\times10^8\times\frac{\sin 143.71^\circ}{\sin 0.09^\circ}\\
&\fbox{$d_{\mathrm{Sun}} = 1.48\times10^{11}\,\mathrm{m}$} & &\hspace{-2em}(1.0)
\end{align*}
\end{description}

\solhead{3.36}{Hypothesized Planet in the Solar System}
% source id: 2025_OBS-PLAN_OP01
\begin{description}
\item[(OP01.1)]
The orbits of the 4 TNOs are elliptical, which can be noticed from careful observations of
their speed. Hence, to estimate lengths of their semi-major axes it is essential to measure
the time period of one complete revolution of each TNO. The time periods of the planets are
\[
\mathrm{O1{:}\ 40\ sec},\quad \mathrm{O2{:}\ 30\ sec},\quad
\mathrm{O3{:}\ 48\ sec},\quad \mathrm{O4{:}\ 34\ sec}
\]
Hence, the correct ascending order is: $\mathrm{O2} < \mathrm{O4} < \mathrm{O1} < \mathrm{O3}$.
% The 24-row credit table for all possible orderings (marking scheme) is omitted.

\item[(OP01.2)]
The orbits of the TNOs can be drawn by marking a few important constellations and noticing
the fact that some prominent or important stars lie along the TNOs' paths. Also the
intersection points of each TNO trajectory with respect to the Earth's orbit (i.e.\ the
Ecliptic) lie near prominent stars.
\ioaafig{2025_OBS-PLAN_OP01-s1.pdf}

Since this map has an angle preserving (Stereoscopic) projection, measuring the angles
directly using a protractor at the intersection points gives the inclination of the orbits
with respect to the Ecliptic.

The measured values of inclination are:
\begin{itemize}
\item $i_1 = 18^\circ$
\item $i_2 = 14^\circ$
\item $i_3 = 20^\circ$
\item $i_4 = 12^\circ$
\end{itemize}
The average inclination $i_{hp} = 16^\circ$.
% The per-TNO trace/angle marking scheme lines are omitted.
\end{description}

\solhead{3.37}{Observer on a Ringed Planet}
% source id: 2025_OBS-PLAN_OP02
\begin{description}
\item[(OP02.1)]
% The solution is the marked sky map: the positions of the 10 moons are indicated by
% scoring circles on the solution overlay.
\ioaafig{2025_OBS-PLAN_OP02-s1.pdf}
For each moon:
\begin{itemize}
\item Centre of the $\otimes$ within the inner circle \hfill 2.0
\item Centre of the $\otimes$ between the inner and outer circles \hfill 1.0
\end{itemize}
Use 4 correctly placed moons for marking. Every incorrectly placed moon beyond 4 gets
$-1.0$, upto a minimum of 0. % sic: "upto"

\item[(OP02.2)]
The different angular extents of the shadow on the inner and outer edges of the ring are
visible on the dome. These edges may be marked on the sky map (using convenient positions
of the stars nearby).

From the grid of the RA lines on the map the angular sizes are measured to be:
\begin{center}
\begin{tabular}{ll}
Inner edge: & $\approx 75^\circ$\\
Outer edge: & $\approx 33^\circ$\\
\end{tabular}
\end{center}
These angles need to be used in the schematic diagram to locate the rings.

Place a protractor with its baseline along the dotted line, and with its origin at P (the
position of the observer).

Find the line of sight to the inner edge that subtends an angle of $75^\circ/2 = 37.5^\circ$
at P. Extend this line to intersect the tangent ray. This point marks one end of the shadow
on the inner ring. Similarly, find the other ends of the shadow on the inner edge.

Draw a circular arc (with centre at the centre of the planet) passing through these two
points of intersection to show the inner edge of the ring. Repeat for the outer edge.
\ioaafig{2025_OBS-PLAN_OP02-s2.pdf}

\item[(OP02.3)]
Once the edges of the ring are drawn, the width can be calculated by measuring the
difference of their distances from the centre of the planet. Use a ruler to measure these
distances. Also measure the radius of the planet on the diagram with the ruler to find the
calibration.

We find,
\begin{itemize}
\item Distance between the two rings along the dotted line: 3.8 cm
\item Radius of the planet on the diagram: 2 cm equivalent to $6.00\times10^4$ km (given)
\end{itemize}
Therefore,
\[
\text{Width of the ring in km:}
= \frac{3.8\,\mathrm{cm}}{2\,\mathrm{cm}}\times 6.00\times10^4\,\mathrm{km}
= 11.4\times10^4\,\mathrm{km}
\]

Tolerance on values of radii of inner and outer edges:
\begin{center}
\begin{tabular}{|c|c|c|}
\hline
Half credit lower limit & Full credit range & Half credit upper limit\\
\hline
$9.45\times10^4$ km & $10.60\times10^4$ km to $13.40\times10^4$ km & $14.55\times10^4$ km\\
\hline
\end{tabular}
\end{center}
\end{description}

\solhead{3.38}{Retrograde Mars}
% source id: 2023_OBS-PLAN_2
% Official solution (2023_PLAN_sol.pdf, Planetarium Round 2 'Retrograde Mars', p.71).
\begin{description}
\item[(a)] 2 points for correct date of opposition and 1 point for each other answer
(total 8)
\begin{center}
\begin{tabular}{|ll|l|}
\hline
i. & the dates of quadrature (when the elongation of Mars is $90^\circ$)
  & 8 Nov 1915 $\pm$ 5 days\\
 & & 15 May 1916 $\pm$ 5 days\\
\hline
ii. & the date of the beginning of retrograde motion & 1 Jan 1916 $\pm$ 5 days\\
 & and the date of the end of retrograde motion & 21 Mar 1916 $\pm$ 5 days\\
\hline
iii. & the date of opposition & 20 Feb 1916 $\pm$ 5 days\\
\hline
iv. & the ecliptic latitude at opposition & $4.5^\circ \pm 1^\circ$\\
\hline
v. & the width in ecliptic longitude of the loop made by the planet
  & $19^\circ \pm 2^\circ$\\
\hline
\end{tabular}
\end{center}

\item[(b)] radius 1.5 au.
% No worked construction is printed; the answer sheet (p.70) shows only the blank
% heliocentric diagram (Earth's orbit with the Sun at centre and the Earth on it) on which
% the geometric construction was to be drawn.

\item[(c)] inclination $1.5^\circ$.
\end{description}
